\documentclass[11pt]{article}
\usepackage[margin=1.05in]{geometry}
\usepackage{amsmath,amssymb,amsthm,mathtools}
\usepackage{microtype}
\usepackage{booktabs}
\usepackage{enumitem}
\usepackage{longtable,array,xurl}
\newcommand{\Lam}{\Lambda}
\newcommand{\one}{\mathbf 1}
\newcommand{\omegaR}{\omega}
\newcommand{\MR}{\mathsf M}
\newcommand{\JR}{\mathcal J_R}
\newcommand{\F}{\mathfrak f}
\newcommand{\R}{\mathbb R}
\newcommand{\epsp}{\epsilon_p}
\newcommand{\csmall}{\mathrm{CORE}}
\newcommand{\clong}{\mathrm{LONG}}
\usepackage[hidelinks,bookmarksopen=true,bookmarksnumbered=true]{hyperref}
\hypersetup{pdftitle={Pointwise Transfer of Anchored Singular Series at Fixed Modulus},pdfauthor={Charlie Dorsey},pdfsubject={Mathematical article with integrated analytic proof appendices},pdfkeywords={singular series, fixed modulus, analytic number theory, formal proof}}

\usepackage[nameinlink,noabbrev]{cleveref}

\newtheorem{theorem}{Theorem}[section]
\newtheorem{proposition}[theorem]{Proposition}
\newtheorem{lemma}[theorem]{Lemma}
\newtheorem{corollary}[theorem]{Corollary}
\theoremstyle{definition}
\newtheorem{definition}[theorem]{Definition}
\theoremstyle{remark}
\newtheorem{remark}[theorem]{Remark}

\newcommand{\SSer}{\mathfrak S}
\newcommand{\Mob}{\mathsf M}
\newcommand{\ZE}{\mathsf Z_E}
\newcommand{\evE}{\operatorname{ev}_E}
\newcommand{\1}{\mathbf 1}
\newcommand{\e}{\mathrm e}
\newcommand{\ph}{\varphi}
\newcommand{\Om}{\Omega}
\newcommand{\Dcal}{\mathcal D}
\newcommand{\Fcal}{\mathcal F}
\newcommand{\Kcal}{\mathcal K}
\newcommand{\Lcal}{\mathcal L}
\newcommand{\Zmod}{\mathbb Z}
\newcommand{\eps}{\varepsilon}

\title{Pointwise Transfer of Anchored Singular Series at Fixed Modulus}
\author{Charlie Dorsey}
\date{October 2026}

\begin{document}
\maketitle

\begin{abstract}
For a fixed positive squarefree modulus $q$ and compatible reduced
endpoint classes $a,b\pmod q$, we establish the pointwise level-one
singular-series transfer
\[
 A_{1,q}(a,b;h)-B_{1,q}(a,b;h)=o_q(\log h),
\]
uniformly over eligible endpoints as the positive compatible height
$h$ tends to infinity. A finite Boolean contraction identity gives
an exact dictionary between completed endpoint coordinates and
connected free-label coordinates. The analytic proof retains its
original signed arithmetic source, partitions it into three
cardinality sectors and five CORE/seven LONG physical regions,
and obtains the logarithmic saving after a single decoration
aggregation. Independent Euler-product identities connect this
analytic estimate to the literal quantities in the theorem.
No density-one height restriction or conjecture on actual
consecutive-prime frequencies is used.
\end{abstract}

\section{Introduction and scope}\label{sec:intro}

For a finite set $H\subset\Zmod$ and a squarefree integer $q\ge1$, define the $q$-restricted singular series
\begin{equation}\label{eq:S_q}
\SSer_q(H)
=\prod_{p\nmid q}
\left(1-\frac{\nu_p(H)}p\right)
\left(1-\frac1p\right)^{-|H|},
\end{equation}
where $\nu_p(H)$ is the number of residue classes occupied by $H$ modulo $p$.  As usual, the product is zero when the tuple is inadmissible at an active prime.  Let
\begin{equation}\label{eq:S_q0}
\SSer_{q,0}(H)
=\sum_{J\subseteq H}(-1)^{|H|-|J|}\SSer_q(J)
\end{equation}
be its Boolean M\"obius transform, in the modified singular-series tradition of Montgomery--Soundararajan \cite{MontgomerySoundararajan2004}.

The central question of this paper is whether the connected
contribution of a single interior position relative to two fixed
endpoints can, after summation, be replaced by its two one-anchor
boundary contributions at the logarithmic scale. These
quantities occur naturally in fixed-gap singular-series models,
but equality at leading scale is an arithmetic assertion about
Euler coefficients, not an assertion about actual prime gaps.
The theorem gives the pointwise answer for every sufficiently
large compatible height at each fixed squarefree conductor.

The first distinction is between an \emph{exact identity} and
an \emph{analytic estimate}. The endpoint contraction formula in
\cref{sec:dictionary} is an identity in a finite Boolean lattice;
it says what a completed connected coordinate is. The exact
Euler bridge in \cref{sec:euler-bridge} says that the quantities
printed in the theorem agree with independently defined arithmetic
kernels, including their infinite-product normalizations. The
analytic argument in \cref{sec:signed-source,sec:cells,sec:aggregation}
then proves a small signed difference of those kernels.
Neither branch can be replaced by the other.

The second distinction is between \emph{preserving a source} and
\emph{estimating it}. The proof starts with a decorated signed sum
whose coefficients depend on a moving reduced height, proper
ordered factorizations and CRT potentials. A sequence of exact
partitions assigns each original summand to a region on which
one available method really applies. The difficult CORE and LONG
sectors are then bounded through twelve disjoint physical cells.
Some obtain oscillatory cancellation; others are sufficiently thin
or arithmetically sparse that direct mass estimates are better.
The full savings survive the outer decoration weights by a strict,
though very small, exponent margin. The optional reader companion, \emph{Reading the Proof},
provides a route from the theorem to this complete article, its
appendices and the precise formal-source waypoints. The technical
appendices make four especially delicate mechanisms independently
inspectable: the CORE band, the LONG edge and its complementary
Weil region, the two physical reciprocal-Suzuki routes, and the
exact B-side centering tail. It states the remaining cell contracts
without suggesting that a bound on one parent is a bound on the
corresponding signed physical cell.

The modified singular-series transform originates with Montgomery--Soundararajan \cite{MontgomerySoundararajan2004}; fixed anchors and the conditional fixed-gap framework have
well-developed precedents in the consecutive-prime literature
\cite{GoldstonLedoan2011,LemkeOliverSound2016}. Restricted
singular-series sums were previously studied, for example, by
Kuperberg \cite{Kuperberg2025}. The present contribution is a
specific \emph{unconditional coefficient transfer at fixed
squarefree conductor}, together with its exact source-sensitive
proof. It is not a claim of priority for CRT, Boolean inversion,
Vaaler approximation, Vinogradov's method or Weil bounds.

\section{The connected contraction dictionary}\label{sec:dictionary}

Fix the endpoint set
\[
E=\{0,h\}.
\]
Let $T$ be any finite set of free labels disjoint from $E$.  For $A\subseteq E$ and $U\subseteq T$, put
\[
f_q(A,U)=\SSer_q(A\cup U).
\]
Define partial Boolean M\"obius operators
\begin{align}
(\Mob_E f)(A,U)
&=\sum_{B\subseteq A}(-1)^{|A|-|B|}f(B,U),\label{eq:ME}\\
(\Mob_T f)(A,U)
&=\sum_{V\subseteq U}(-1)^{|U|-|V|}f(A,V).\label{eq:MT}
\end{align}
They commute, and
\[
(\Mob_E\Mob_Tf_q)(A,U)=\SSer_{q,0}(A\cup U).
\]
Let endpoint-fiber contraction be
\[
(\ZE g)(U)=\sum_{A\subseteq E}g(A,U),
\]
and let $\evE$ denote evaluation at the full endpoint set $A=E$.

\begin{proposition}[Endpoint contraction identity]\label{prop:endpoint-contraction}
For every $q$, every endpoint pair $E$, and every finite free set $T$,
\begin{equation}\label{eq:commuting-identity}
\boxed{\ZE\Mob_E\Mob_T=\Mob_T\evE.}
\end{equation}
Equivalently,
\begin{equation}\label{eq:F_E}
\Fcal_{E,q}(T)
:=\sum_{A\subseteq E}\SSer_{q,0}(A\cup T)
=\Mob_T[\SSer_q(E\cup\cdot)](T).
\end{equation}
\end{proposition}

\begin{proof}
For fixed $U\subseteq T$,
\[
\sum_{A\subseteq E}\sum_{B\subseteq A}
(-1)^{|A|-|B|}f_q(B,U)
=f_q(E,U),
\]
which is the Boolean zeta--M\"obius identity on the endpoint fiber.  Since $\Mob_E$ and $\Mob_T$ commute, \eqref{eq:commuting-identity} follows.
\end{proof}

Assume $\SSer_q(E)>0$ and define
\begin{equation}\label{eq:Rq}
R_q(U)=\frac{\SSer_q(E\cup U)}{\SSer_q(E)},
\qquad
R_q^\circ(T)=\sum_{V\subseteq T}(-1)^{|T|-|V|}R_q(V).
\end{equation}
Then \cref{prop:endpoint-contraction} gives the pointwise identity
\begin{equation}\label{eq:F-Rcirc}
\boxed{\Fcal_{E,q}(T)=\SSer_q(E)R_q^\circ(T).}
\end{equation}
If $I$ is a finite interior set and
\[
C_{k,q}^{\mathrm{comp}}
=\sum_{\substack{T\subseteq I\\|T|=k}}\Fcal_{E,q}(T),
\qquad
A_{k,q}^{\circ}
=\sum_{\substack{T\subseteq I\\|T|=k}}R_q^\circ(T),
\]
then
\begin{equation}\label{eq:C-A-dictionary}
\boxed{C_{k,q}^{\mathrm{comp}}=\SSer_q(E)A_{k,q}^{\circ}.}
\end{equation}
At $q=1$ this is the exact normalization dictionary between the completed endpoint coordinate and the anchored free-label connected coordinate.

\subsection{Anchor subtraction}

Write
\[
c_q(A,T)=\SSer_{q,0}(A\cup T),
\]
and define the lower-anchor sum and the top-anchor sector by
\begin{align}
b_q(T)&=c_q(\varnothing,T)+c_q(\{0\},T)+c_q(\{h\},T),\label{eq:bq}\\
g_q(T)&=c_q(E,T)=\SSer_{q,0}(E\cup T).\label{eq:gq}
\end{align}
Then $\Fcal_{E,q}(T)=b_q(T)+g_q(T)$.  Since singleton singular series are one,
\[
\SSer_{q,0}(E)=\SSer_q(E)-1.
\]
Combining this with \eqref{eq:F-Rcirc} gives the exact residual identity
\begin{equation}\label{eq:anchor-residual}
\boxed{
R_q^\circ(T)-b_q(T)
=\frac{g_q(T)-\SSer_{q,0}(E)b_q(T)}{\SSer_q(E)}.
}
\end{equation}
Thus the anchored-to-free transfer defect is the normalized failure of the top two-anchor connected sector to contract onto the lower anchor sectors.

\subsection{The same transform at one prime}

At an active prime $p$, let $z$ be uniformly distributed modulo $p$ and
introduce the \emph{mean-one avoidance factor}
\begin{equation}\label{eq:local-avoidance}
 \phi_{p,z}(x)=\frac{1-\mathbf 1_{z\equiv -x\pmod p}}{1-1/p},
 \qquad \frac1p\sum_{z\bmod p}\phi_{p,z}(x)=1.
\end{equation}
For a finite set $H$ (and, when coordinates are independently labeled,
for a labeled list of occurrences), the local factor is exactly
\[
 L_p(H)=\frac1p\sum_{z\bmod p}\prod_{x\in H}\phi_{p,z}(x)
       =\frac{1-\nu_p(H)/p}{(1-1/p)^{|H|}}.
\]
In the labeled case, $|H|$ counts coordinate occurrences even when values
repeat; the product retains repeated factors although the numerator
counts occupied residue classes. Boolean connectedization is thus
\begin{equation}\label{eq:local-connected}
 \sum_{J\subseteq H}(-1)^{|H|-|J|}L_p(J)
    =\frac1p\sum_{z\bmod p}\prod_{x\in H}(\phi_{p,z}(x)-1).
\end{equation}
The centering in \eqref{eq:local-avoidance} removes the constant Fourier
coefficient in each \emph{individual} free-coordinate factor. This
single-prime calculation is an algebraic identity, not an estimate on
the entire signed arithmetic sum.

\section{Fixed squarefree modulus and the level-one theorem}\label{sec:fixed-q-setup}

Fix squarefree $q\ge1$ and endpoint classes
\[
a,b\in(\Zmod/q\Zmod)^\times,
\qquad h\equiv b-a\pmod q.
\]
Assume the endpoint pair is $q$-restricted admissible,
\[
\SSer_q(E)>0,
\qquad E=\{0,h\}.
\]
Every legal height is even: if $2\nmid q$, admissibility at the active prime $2$ forces $h$ even; if $2\mid q$, the unit classes $a,b$ are both odd and hence their difference is even modulo $2$, so $h\equiv b-a\equiv0\pmod2$.

Define the admissible interior and its exact size by
\begin{equation}\label{eq:IqNq}
I_q(a;h)=\{0<t<h:(a+t,q)=1\},
\qquad N_q(a;h)=|I_q(a;h)|.
\end{equation}
The level-one connected and lower-anchor sums are
\begin{align}
A_{1,q}(a,b;h)
&=\sum_{t\in I_q(a;h)}R_q^\circ(\{t\})\notag\\
&=\sum_{t\in I_q(a;h)}
\left(
\frac{\SSer_q(\{0,t,h\})}{\SSer_q(\{0,h\})}-1
\right),\label{eq:A1q}\\
B_{1,q}(a,b;h)
&=\sum_{t\in I_q(a;h)}b_q(\{t\})\notag\\
&=\sum_{t\in I_q(a;h)}
\left[\SSer_{q,0}(\{t\})
+\SSer_{q,0}(\{0,t\})
+\SSer_{q,0}(\{h,t\})\right].\label{eq:B1q}
\end{align}
Since $\SSer_{q,0}(\{t\})=0$,
\begin{equation}\label{eq:B1q-simple}
B_{1,q}
=\sum_{t\in I_q(a;h)}(\SSer_q(\{0,t\})-1)
+\sum_{t\in I_q(a;h)}(\SSer_q(\{h,t\})-1).
\end{equation}
By \eqref{eq:anchor-residual},
\begin{equation}\label{eq:AminusB-connected}
A_{1,q}-B_{1,q}
=\frac1{\SSer_q(E)}
\sum_{t\in I_q(a;h)}
\left[g_q(\{t\})-\SSer_{q,0}(E)b_q(\{t\})\right].
\end{equation}

\begin{theorem}[Fixed-squarefree-$q$ pointwise level-one transfer]\label{thm:fixed-q}
Fix a positive squarefree integer $q$. For every $\varepsilon>0$
there exists $h_0=h_0(q,\varepsilon)$ such that every positive
$h\ge h_0$ and every reduced pair of endpoints $a,b\pmod q$
with $a+h\equiv b\pmod q$ and
$2\nmid q\Rightarrow 2\mid h$ satisfy
$|A_{1,q}(a,b;h)-B_{1,q}(a,b;h)|\le\varepsilon\log h$.
Equivalently, uniformly over those compatible endpoints,
\begin{equation}\label{eq:main-theorem}
\boxed{A_{1,q}(a,b;h)-B_{1,q}(a,b;h)=o_q(\log h)}
\end{equation}
as $h\to\infty$ through positive compatible heights.
\end{theorem}

There is no density-one restriction in $h$.

\subsection{Exact density normalization}

The admissible count is exactly periodic.

\begin{lemma}[Exact periodic count]\label{lem:Nq}
For each compatible pair $(a,b)$,
\begin{equation}\label{eq:Nq-periodic}
\boxed{N_q(a;h)=\frac{\ph(q)}q h+\eps_q(a,b),}
\end{equation}
where $\eps_q(a,b)$ depends only on the fixed compatible residue data, not on the positive representative of the gap class.
\end{lemma}

\begin{proof}
If $h$ is increased by $q$ while preserving its residue class, the interval gains one complete period of $q$ translates and hence exactly $\ph(q)$ admissible interior positions.  Therefore $N_q(a;h)-\ph(q)h/q$ is constant on each compatible height progression.
\end{proof}

For the completed Lemke Oliver--Soundararajan normalization, the fixed-modulus density parameter is
\begin{equation}\label{eq:deltaq}
\delta_q(x)=\frac{q}{\ph(q)\log x}.
\end{equation}
Consequently
\[
\delta_q(x)\frac{\ph(q)}q=\frac1{\log x}.
\]
The $q/\ph(q)$ factor belongs to the ambient admissible-point density; it does not appear as an additional multiplier inside the transfer kernel \eqref{eq:main-theorem}.  The unconditional level-one contraction formula is
\begin{equation}\label{eq:LOS-level-one}
 C_{1,q}^{\mathrm{comp}}=\SSer_q(E)A_{1,q}.
\end{equation}
Only when $N_q(a;h)>0$ is the completed coefficient per
admissible interior position defined by
\begin{equation}\label{eq:LOS-level-one-per-site}
 \Delta\mu_0:=\frac{C_{1,q}^{\mathrm{comp}}}{N_q(a;h)}
 =\frac{\SSer_q(E)}{N_q(a;h)}A_{1,q}.
\end{equation}
This qualification is necessary: for $q=2$, $a=b=1$, $h=2$
the only interior integer is not admissible, so $N_q(a;h)=0$.
By \cref{lem:Nq}, however, for each fixed compatible residue pair
$N_q(a;h)=(\ph(q)/q)h+\eps_q(a,b)>0$ eventually. Therefore the
finite-height domain restriction does not affect
\cref{thm:fixed-q}.

\section{Exact representation of the Euler quantities}\label{sec:euler-bridge}

The quantities in \eqref{eq:A1q} and \eqref{eq:B1q} are finite sums of
\emph{infinite} products. Their analytic treatment starts with an exact
identification, not an asymptotic replacement. Write
$A^{\mathrm{src}}_{1,q}(a;h)$ and $B^{\mathrm{src}}_{1,q}(a,b;h)$ for the
arithmetic divisor-periodic kernels defined in the accompanying source
(see the source-obligation concordance distributed with the accompanying analytic supplement and the reader companion). We deliberately keep the
superscript: equal numerical values do not make the two descriptions
interchangeable without proving that the normalizations agree.

\begin{proposition}[Literal Euler-to-source identity]\label{prop:euler-bridge}
Let $q\ge1$, let $a+h\equiv b\pmod q$, and impose the parity condition
$2\nmid q\Rightarrow2\mid h$. For the literal products in
\eqref{eq:S_q}, \eqref{eq:A1q}, and \eqref{eq:B1q},
\begin{equation}\label{eq:euler-source-bridge}
 A_{1,q}(a,b;h)-B_{1,q}(a,b;h)
 =A^{\mathrm{src}}_{1,q}(a;h)-B^{\mathrm{src}}_{1,q}(a,b;h).
\end{equation}
The identity retains the exceptional active primes, the inactive factors
for $p\mid q$, and the complete centering of the boundary terms.
\end{proposition}

\begin{proof}
We separate two points that are easy to conflate. For the anchored
quotient, primewise division of the three-point and endpoint-pair
factors requires the \emph{whole} endpoint product to be nonzero.
The high-prime endpoint factors are of the form
\[
 1-\frac{1}{(p-1)^2}
\]
when the two endpoints occupy distinct classes; their deviations from
$1$ are summable. The finitely many exceptional factors, including the
active prime $2$, are handled according to the parity/admissibility
condition. Only after convergence and nonvanishing are established is
it legitimate to identify the quotient of products with the product of
local quotients. Each active factor then has an exact divisor-selector
expansion. When summing those selectors over the admissible interior,
the subtraction is the \emph{exact} cardinality $N_q(a;h)$, not a
mean-density approximation.

The boundary expression requires a different argument. At an active
high prime the free-pair factor has the selector
\[
 c_1(p)\bigg(1+\frac{\mathbf1_{p\mid y}}{p-2}\bigg),
 \qquad c_1(p)=\frac{p(p-2)}{(p-1)^2}.
\]
More explicitly, the high-prime divisor coefficient and constant are
\[
 \eta_q(e)=\mu^2(e)\mathbf1_{(e,6q)=1}
       \prod_{p\mid e}\frac1{p-2},\qquad
 C_{f,q}=\prod_{\substack{p\ge5\\p\nmid q}}
       \frac{p(p-2)}{(p-1)^2}.
\]
All primes in the coefficient support are at least $5$ and do not
divide $q$. The exact Euler normalization then gives
\begin{equation}\label{eq:boundary-normalization}
 C_{f,q}\sum_{e\ge1}\frac{\eta_q(e)}{e}=1.
\end{equation}
The complete first-anchor expression also carries the exceptional active
primes $2$ and $3$. For a residue representative $r\pmod q$ define
\begin{equation}\label{eq:free-low-weight}
 w_{q,r}(t)=\mathbf1_{(r+t,q)=1}
 \prod_{\substack{p\in\{2,3\}\\p\nmid q}}
 \frac{1-\nu_p(\{0,t\})/p}{(1-1/p)^2},
 \qquad
 R^{w}_e(q,r;h)=\sum_{\substack{0<t<h\\e\mid t}}w_{q,r}(t).
\end{equation}
The exact centered \emph{first-anchor} identity is
\begin{equation}\label{eq:boundary-infinite-tail}
 \boxed{S_1(q,r;h)=C_{f,q}\sum_{e\ge1}\eta_q(e)
     \left(R^w_e(q,r;h)-\frac{N_q(r;h)}e\right).}
\end{equation}
Here $S_1(q,r;h)=\sum_{0<t<h,\,(r+t,q)=1}
(\SSer_q(\{0,t\})-1)$ is the literal first-boundary contribution.
It is crucial that the \emph{positive mass is weighted but the
negative centering is not}: replacing $N_q(r;h)/e$ by a weighted
mean would give a different normalization. Equation
\eqref{eq:boundary-normalization} supplies the full subtraction
$-N_q(r;h)$, and the convergence argument justifies passage between
the finite divisor-swapped numerator and the infinite centered series.
The numerator vanishes for $e\ge h$; the negative tail does not.

For a simple check, take $q=1,h=2$. The only interior displacement
$t=1$ is annihilated by the active prime $2$, so the literal
first-boundary sum is $-1$. Correspondingly $w_{1,0}(1)=0$ and
\eqref{eq:boundary-normalization} makes
\eqref{eq:boundary-infinite-tail} equal to $-1$. An unweighted
progression count would instead give $C_{f,1}-1$, which is different.

The second boundary must be treated by a separate reflection. Let
$r_b$ represent $-b\pmod q$. Under $t\mapsto h-t$ and
$a+h\equiv b\pmod q$, the conditions
$(a+t,q)=1$ and $(r_b+h-t,q)=1$ agree. Translation of each
restricted pair Euler factor then gives the \emph{exact} identity
\begin{equation}\label{eq:two-boundary-weighted}
 \boxed{B_{1,q}(a,b;h)=S_1(q,a;h)+S_1(q,r_b;h),}
\end{equation}
with each $S_1$ represented by \eqref{eq:boundary-infinite-tail} and
its own (unweighted) admissible count. The exact reflection also
identifies those two admissible counts on compatible heights. The
A-side and both B-side Euler bridges, including their convergence
and translation steps, establish \eqref{eq:euler-source-bridge}.
\end{proof}

It remains to estimate the right-hand side of
\eqref{eq:euler-source-bridge}. The analytic proof does \emph{not}
rederive \eqref{eq:euler-source-bridge} as an approximation; the two
proof obligations remain separate until the final theorem.

\section{The signed physical source and its exact partitions}
\label{sec:signed-source}

\subsection{The original finite signed source}

Put $\Lambda=\log h$ and $\sigma_0=\frac23+\frac1{100}$, and define
the actual source cutoff
\begin{equation}\label{eq:source-Z}
 Z(h)=2^{20(\lfloor\Lambda^{\sigma_0}\rfloor+1)}.
\end{equation}
A decoration $d$ is a subset of the prime divisors of $h$ which
are at least $5$ and do not divide $q$. Write
$D_d=\prod_{p\in d}p$, $H_d=h/D_d$, and
$V_h(d)=\prod_{p\in d}(p-2)^{-1}$. The decorated sum retains only
those $d$ for which
\begin{equation}\label{eq:active-decoration}
  Z(h)\exp(\Lambda^{1/4})<H_d.
\end{equation}
The corresponding finite high-prime support is
\[
 T_{q,h,d}=\{p\le H_d^2:p\text{ prime},\ p\ge5,\ p\nmid qH_d,
                 \ p\notin d\}.
\]
The original quarter-parent family consists of sets of those primes,
with its strict lower and weak upper product bounds:
\begin{equation}\label{eq:quarter-parent-contract}
 \mathcal Q_{q,h,d}=
  \Bigl\{e\subseteq T_{q,h,d}:
   Z(h)<\prod_{p\in e}p,\quad
   4\prod_{p\in e}p\le H_d^2\Bigr\}.
\end{equation}
For each such $e$ we retain every \emph{proper, nonempty,
ordered} split $e=B\sqcup C$, with
$r=\prod_{p\in B}p$, $s=\prod_{p\in C}p$.
The Central gate is exactly
\begin{equation}\label{eq:central-contract}
 rs\le H_d e^{-\Lambda^{1/4}},\qquad
 s\le r^2,\qquad r\le s^2.
\end{equation}
The parent weight is $\omega(e)=\prod_{p\in e}(p-3)^{-1}$.
The first condition in \eqref{eq:central-contract} is a genuine
restriction on the ordered pair, not an optional asymptotic estimate.

To specify the \emph{signed} row, let $a_d\equiv D_d^{-1}a\pmod q$
(the inverse exists), and define the exceptional-prime occupied weight
\[
 \lambda_{q,a_d,H}(y)=\mathbf1_{(a_d+y,q)=1}
 \prod_{\substack{p\in\{2,3\}\\p\nmid q}}
 \frac{1-\nu_p(\{0,y,H\})/p}{1-\nu_p(\{0,H\})/p}
 (1-1/p)^{-1}.
\]
The mean and centered weight over its source period $6q$ are
\[
 \mu_{q,a_d,H}=\frac1{6q}\sum_{y=0}^{6q-1}\lambda_{q,a_d,H}(y),
 \qquad g_{q,a_d,H}(y)=\lambda_{q,a_d,H}(y)-\mu_{q,a_d,H}.
\]
For $n$ coprime to $M=6q$ define the periodic arithmetic potential
\[
 \mathcal P_{g,M,n}(y)=\sum_{j=0}^{[yn^{-1}]_M-1}g(jn),
\]
where $[\cdot]_M$ is the representative in $\{0,\ldots,M-1\}$.
Let $\bar r_s$ be the inverse of $r$ modulo $s$ (the
permitted source pairs are coprime), and define the CRT residue
\[
 c_{H;r,s}=r\,[H\bar r_s]_s,\qquad 0\le [H\bar r_s]_s<s,
\]
where $[\cdot]_s$ denotes the least nonnegative residue.  This
is the source's literal \texttt{crtResidue}$(H,r,s)$, not a
choice of arbitrary CRT representative.  With the
source sawtooth $\psi(x)=\{x\}-\frac12$, including its
integer convention, the folded row is
\begin{equation}\label{eq:folded-row-contract}
\begin{split}
 \mathfrak f_{q,a,D_d}(H_d,r,s)
   ={}&-2\mu_{q,a_d,H_d}\,
      \psi\!\left(\frac{H_d\bar r_s}s\right)\\
     &+\mathcal P_{g,6q,rs}(H_d+rs-c_{H_d;r,s})
       -\mathcal P_{g,6q,rs}(c_{H_d;r,s}).
\end{split}
\end{equation}
This expression is the exact folded combination of the oriented
reciprocal sawtooth and the two CRT-potential evaluations. The
mean subtraction in $g$ and the ordering of the two potentials
are part of the analytic data.

Finally the \emph{complete signed folded Central sum} is
\begin{equation}\label{eq:physical-source-schema}
\boxed{
 \mathcal C_{\rm fold}(q,a,h)=
 \sum_{\substack{d\subseteq\mathcal P_{q,h}\\Z(h)e^{\Lambda^{1/4}}<H_d}}
 V_h(d)\!
 \sum_{e\in\mathcal Q_{q,h,d}}
 \sum_{\varnothing\ne B\subsetneq e}
 \mathbf1_{\mathsf{Central}(H_d,r,s)}
 \omega(e)\,\mathfrak f_{q,a,D_d}(H_d,r,s),
}
\end{equation}
where $\mathcal P_{q,h}=\{p\mid h:p\text{ prime},p\ge5,p\nmid q\}$,
$r=\prod_{p\in B}p$, and $s=\prod_{p\in e\setminus B}p$.
This is the particular finite signed source to which the cell
estimates apply; none is inferred for an arbitrary row function.
The corresponding source definition initially ranges over all legal
decoration subsets. Its lower-product gate forces the active condition
\eqref{eq:active-decoration}; hence every inactive decoration term
vanishes, proving that the active-only display and the all-decoration
source are exactly equal. This is an identity of sums, not an
asymptotic deletion of small terms.

The original active Central source and folded Central have the
independently proved exact identity
\begin{equation}\label{eq:active-folded}
 \mathcal C_{\rm active}=
  \mathcal C_{\rm fold}+\mathcal R_{\rm inv}+\mathcal R_{\rm dens}.
\end{equation}
The error terms are separately bounded in \cref{sec:aggregation}.
The source contract and its identity are recorded by the
physical source definition, the signed-fold theorem, and the
Central transfer. The present paper uses these \emph{exact}
representations, not a schematic surrogate for a difficult sum.

\subsection{An identity, not a cancellation estimate}

Choose cutoff functions $J_\ell(h)\le J_u(h)$. Every finite
parent has exactly one of the three cardinalities
\[
 |e|<J_\ell(h),\qquad
 J_\ell(h)\le|e|<J_u(h),\qquad
 J_u(h)\le|e|.
\]
Multiplying \eqref{eq:physical-source-schema} by the corresponding
indicators and summing gives the exact signed equality
\begin{equation}\label{eq:sector-partition}
 \boxed{\mathcal C_{\rm fold}
  =\mathcal C_{\rm LOW}+\mathcal C_{\rm MIDDLE}
   +\mathcal C_{\rm HIGH}.}
\end{equation}
No absolute value, summation inequality, or large-height assumption
has yet been used. This matters: the original signs, decorations and
ordered pairs must remain attached to the summands on both sides.

The middle population is split by the geometric test
\begin{equation}\label{eq:long-guard}
 \mathcal L_\delta(r,s):
 \quad s^{1+\delta}<r\quad\text{or}\quad r^{1+\delta}<s,
 \qquad 0<\delta<\frac13.
\end{equation}
CORE is the complement of \eqref{eq:long-guard} \emph{inside the
same middle population}; LONG is its truth set. The two LONG
orientations cannot both hold for positive $r,s$ and $\delta$.
Hence a second exact equality holds:
\begin{equation}\label{eq:middle-partition}
 \boxed{\mathcal C_{\rm MIDDLE}
   =\mathcal C_{\rm CORE}+\mathcal C_{\rm LONG}.}
\end{equation}
The proof is an indicator partition of every original summand,
including the boundary case of equality in the strict tests.

\begin{remark}[Why the lower cutoff may be zero]\label{rem:zero-low}
The analytic bounds on CORE and LONG hold for \emph{arbitrary}
cutoff functions. We therefore set $J_\ell\equiv0$. As
$|e|\ge0$ for each finite parent, the LOW sum is empty and
\begin{equation}\label{eq:low-zero}
 \mathcal C_{\rm LOW}=0,\qquad
 \mathcal C_{\rm fold}
 =\mathcal C_{\rm MIDDLE}(0,J_u)+\mathcal C_{\rm HIGH}(J_u).
\end{equation}
This does \emph{not} assert that the parents previously assigned to
LOW have zero weight. With the new partition their terms belong
to MIDDLE, where the cutoff-uniform bounds apply. Equation
\eqref{eq:sector-partition} is the conservation statement that
makes this reallocation legitimate. We take $\delta=1/4$ at the
final specialization.
\end{remark}

\section{Twelve physical cells and the methods they require}
\label{sec:cells}

The signed source is now partitioned a third time. The five CORE
cells and seven LONG cells are masks on
\eqref{eq:physical-source-schema}, not twelve unrelated
model sums. They are disjoint and exhaustive before estimates.

\begin{proposition}[Exact five- and seven-cell identities]
\label{prop:physical-decomposition}
At each allowed decoration $d$ and at every eligible choice of
cutoffs, one has the exact signed identities
\begin{align}
 \mathcal C_{{\rm CORE},d}
 &=W_d+C_d+R_d+\mathrm{ResA}_d+\mathrm{Band}_d,
 \label{eq:core-five}\\
 \mathcal C_{{\rm LONG},d}
 &=WL_d+CL_d+RL_d+\mathrm{ManyL}_d
    +\mathrm{EdgeL}_d+\mathrm{ResL}_d+WRL_d.
 \label{eq:long-seven}
\end{align}
The original signed folded row, ordered orientation, source gate
and decoration weight are unchanged in both identities.
\end{proposition}

\begin{proof}
CORE tests successively for a prime packet on its first factor,
for the complementary Parseval/Vaughan regime, and for a usable
reciprocal-progression packet on its second factor. Its remaining
parents fall into a rough residual or a near-diagonal band. LONG
first distinguishes the orientation in \eqref{eq:long-guard}. The
longer-left orientation splits into the packet and complementary
Parseval cells. On the longer-right orientation, the tests proceed
through a reciprocal packet, excessive prime-factor count, the
square-root edge, and finally the presence or absence of a Weil
packet. Each test is accompanied by its logical complement;
therefore every admitted parent enters exactly one cell. At the
source level these are the finite \texttt{split5} and
\texttt{split7} identities, applied termwise before summation.
\end{proof}

\begin{table}[htbp]
\centering\small
\caption{CORE: the dominant positive exponent comes from the geometric band.}
\label{tab:core-cells}
\begin{tabular}{@{}lp{3.95in}r@{}}
\toprule
Cell & Applicable source mechanism & Eventual bound \\
\midrule
$W$ & First-factor prime packet; Ramar\'e extraction \cite{Ramare2010}, coupled Cauchy and CRT/Kloosterman & $\Lambda^{-80}$\\
$C$ & Packet absent; Parseval and Vaughan/prime-sum bounds, with Rankin tail & $\Lambda^{-15}$\\
$R$ & Reciprocal-progression packet; shifted Suzuki estimate and Abel summation & $\Lambda^{-75}$\\
$\mathrm{ResA}$ & Neither packet; direct mass of two rough factors & $\Lambda^{1/10}$\\
$\mathrm{Band}$ & Narrow logarithmic diagonal strip; weighted interval mass & $\Lambda^{499/500}$\\
\bottomrule
\end{tabular}
\end{table}

\begin{table}[htbp]
\centering\small
\caption{LONG: geometric control and arithmetic cancellation have separate domains.}
\label{tab:long-cells}
\begin{tabular}{@{}lp{3.95in}r@{}}
\toprule
Cell & Applicable source mechanism & Eventual bound \\
\midrule
$WL$ & Longer-left orientation, packet extraction with Kloosterman control & $\Lambda^{-80}$\\
$CL$ & Longer-left complement, Parseval/Vaughan and exceptional tail & $\Lambda^{-15}$\\
$RL$ & Longer-right reciprocal packet, source-compatible Suzuki input & $\Lambda^{-75}$\\
$\mathrm{ManyL}$ & Many prime factors, suppressed by a Rankin bound & $\Lambda^{-50}$\\
$\mathrm{EdgeL}$ & Short cofactor at square-root edge, bounded by geometric mass & $\Lambda^{19/20}$\\
$\mathrm{ResL}$ & Off edge without a Weil packet; rough-factor mass & $\Lambda^{7/10}$\\
$WRL$ & Off edge with a Weil packet; general Kloosterman, gcd and Cauchy costs & $\Lambda^{-80}$\\
\bottomrule
\end{tabular}
\end{table}

Here and below $\Lambda=\log h$. Each bound controls the
\emph{complete signed masked source sum for a fixed decoration},
uniformly in legal anchors and cutoffs at sufficiently large $h$.
It does not bound the absolute value of every parent separately:
doing so would erase the cancellation that the oscillatory cells
are built to recover. The displayed bounds are the source's \emph{physical} estimates,
with the common quarter-parent, ordered factorization, Central
and cardinality gates inherited throughout. The accompanying
analytic supplement supplies worked geometric and quantitative
chains for the limiting cells, the full remaining cell-contract
register and exact source anchors. That appendix is part of the
mathematical presentation, not a newly claimed independent
reproof of all deep analytic leaf estimates.

\subsection{The physical leaf as a mathematical contract}
\label{sec:leaf-contract}

The operation inside a physical cell has three distinct stages.
First an exact predicate selects terms from the \emph{original}
signed double sum \eqref{eq:physical-source-schema}. Second the
available analytic supplier controls either the resulting
weighted mass or the oscillation of the \emph{whole masked sum};
it is applied only when the packet, cofactor, frequency and
coprimality hypotheses of that region are verified. Third,
constants and any source outer-factor loss are absorbed into
the displayed power of $\Lambda$.

This order matters. In the near-diagonal CORE Band, the two
negations of the earlier geometric tests imply
\begin{equation}\label{eq:band-interval-paper}
 r^{1/(1+\epsilon_p)}<s<r^{1+3\epsilon_p},
 \qquad\epsilon_p=\Lambda^{-1/200}.
\end{equation}
The interval's logarithmic width is at most
$1+4\Lambda^{199/200}$, before a separate roughness bound
on the first factor. In LONG's EdgeL, by contrast, one
\emph{fixes the second factor} and uses
$s^{1/2}\le r<s^{1/2+3\eta}$, $\eta=\Lambda^{-1/10}$,
only after excluding the reciprocal packet and large-factor-count
cases. The geometric operation is not interchangeable with
Weil completion. These two chains, including the floors and
exact roughness/interval constants, are developed in the
supplement rather than buried in the tables above.

\subsection{The exact LONG-right decision chain}

The most revealing comparison is between $\mathrm{EdgeL}$ and $WRL$.
Their support is determined by a sequence of \emph{predicates} on
the same source pairs, not just by a numerical inequality in $r,s$.
Let $T=T_{q,h,d}$, let $J_R(h,T,B)$ be the actual reciprocal-progression
prime window, put
\[
 K_0(h)=\lceil100\log\Lambda\rceil,\qquad
 \eta=\Lambda^{-1/10},
\]
and let $X_p(h)$ and $Y_L(h,s)$ be the source's fixed
packet-window endpoints. On LONG-right the following routing is exact:
\begin{enumerate}[leftmargin=*,itemsep=0.12em]
\item If $C\cap J_R(h,T,B)\ne\varnothing$, assign the pair to $RL$.
\item Otherwise, if $|C|\ge K_0(h)$, assign it to $\mathrm{ManyL}$.
\item Otherwise, if $r<s^{1/2+3\eta}$, assign it to $\mathrm{EdgeL}$.
\item Otherwise, if $B$ meets the source packet
 $\{p\in T\setminus C:X_p(h)\le p\le Y_L(h,s)\}$,
 assign it to $WRL$; if not, assign it to $\mathrm{ResL}$.
\end{enumerate}
These tests inherit the source Central gate, proper ordered split,
cardinality mask, and LONG-right condition; none of those is replaced
by the four new predicates. Thus the short-cofactor strip
\begin{equation}\label{eq:edge-strip}
 s^{1/2}\le r<s^{1/2+3\eta}
\end{equation}
is the \emph{geometric consequence} of the EdgeL mask, not its
entire definition.

Within $\mathrm{EdgeL}$, the cofactor is too short for the
Weil--Cauchy method. The proof uses the weighted logarithmic
width $O(\eta\log s)$ and the roughness forced by missing the
reciprocal packet. Its local cost is exactly
\[
 \underbrace{\tfrac9{10}}_{\rm width}
 +\underbrace{\tfrac1{4000}+\tfrac1{40}+\tfrac1{100}}_{\rm roughness/weight}
 =\tfrac{3741}{4000}<\tfrac{19}{20}.
\]
The final gap $59/4000$ absorbs a fixed coefficient; no Weil estimate
is invoked in this region.

On $WRL$ the prior gates ensure $r\ge s^{1/2+3\eta}$.
Extracting the prime $p\le s^\eta$ leaves a cofactor of length
at least on the scale $s^{1/2+2\eta}$. The available
Kloosterman/Cauchy calculation now applies. Its complete-sum input is
\begin{equation}\label{eq:kloosterman-bound}
 |S(u,v;s)|\le\tau(s)\sqrt{s\,(u,v,s)},
\end{equation}
with its divisor factor and gcd penalty retained
\cite{Weil1948,Estermann1961,IwaniecKowalski2004}.
The coefficient-order reversal in the formal general-modulus
normalization is an exact identity, not a hypothesis that the gcd
factor may be discarded. The source's fixed-factor power
$\Lambda^{-95}$ pays an outer $\Lambda^{10}$ mass, and the
remaining margin absorbs the fixed coefficient to reach the physical
$\Lambda^{-80}$ bound. If the packet is absent, the pair belongs to
$\mathrm{ResL}$, not to $WRL$.

This is the useful geometric principle: cancellation is used only
where the factorization leaves enough range; a sufficiently thin
failure region can be paid by geometric mass instead.


The actual clipped reciprocal application underlying $R$ and $RL$
is worth distinguishing from its final physical estimate. For fixed
$Q=6qr$ and a residue $\rho$, the contributing progression is
restricted \emph{simultaneously} by a dyadic interval and by both
products $vu$ and $v'u$. The resulting natural-number endpoints
$a_*,b_*$ determine $N=(a_*-\rho)\mathbin{\mathrm{div}}Q$ and
$M=(b_*-\rho)\mathbin{\mathrm{div}}Q$. The source proves
$1\le N<M\le2N$, $|Y/v-Y/v'|\ge2Q$ and an Abel
variation norm at most $2$ for its nonconstant cotangent
coefficient. The four signed main/class Fourier channels are
checked separately, under the distinct CORE $5$ and LONG $4$
logarithmic hypotheses. The complete clipped application lemma
and those four coefficient contracts are given in the accompanying
mathematical supplement. Its conclusion is an input to the
source's deeper fixed-factor estimates, not a proof of those
estimates by itself.

\subsection{Why the reciprocal estimate has two versions}

The $R$ and $RL$ cells both use a reciprocal phase of the type
\[
 \sum_{N<m\le U}\exp\!\left(2\pi i\frac{Y}{m+\vartheta}\right),
 \qquad U\le2N.
\]
The source's shifted Suzuki/Karatsuba estimate has a bound of
the shape
\begin{equation}\label{eq:reciprocal-sum-bound}
 \left|\sum_{N<m\le U}e\!\left(\frac{Y}{m+\vartheta}\right)\right|
 \ll N\exp\!\left(-c\frac{(\log N)^3}{(\log Y)^2}\right)
   +\frac{N^2}{Y},
\end{equation}
where $e(t)=\exp(2\pi i t)$ and the fixed constants and
admissible parameter ranges are those of the source theorem.
A low-frequency regime uses a telescoping derivative argument;
a high-frequency regime uses selected derivative orders,
Karatsuba's polynomialization and Vinogradov-type mean-value
bounds. The actual $R$ and $RL$ applications additionally
require progression reindexing, negative-frequency conjugation
and Abel summation for the Fourier coefficients.

The CORE $R$ estimate needs a lower logarithmic factor with
constant $5$. The LONG-right geometry supplies only $4$.
The LONG proof therefore proves the numerical implication
again instead of citing a stronger hypothesis it does not have.
Writing $\sigma_0=\tfrac23+\tfrac1{100}$, the exponent that
matters is
\begin{equation}\label{eq:R-RL-margin}
 3\sigma_0-\frac3{200}-2=\frac3{200}>0.
\end{equation}
The change from $5^3$ to $4^3$ affects the fixed coefficient
but does not destroy this positive power. It is an instance of
transporting a classical mechanism with its genuine quantitative
price rather than treating theorem names as universally portable.

\section{The logarithmic budget and the uniform terminal estimate}
\label{sec:aggregation}

\subsection{The largest cell controls the sector}

\begin{proposition}[Physical source bounds]\label{prop:physical-rates}
For each fixed $q>0$ and $0<\delta<1/3$, there exists
$\Lambda_0(q,\delta)$ such that, whenever
$\Lambda=\log h\ge\Lambda_0$, every permitted decoration $d$,
every legal anchor, and every allowed cardinality cutoff obey
\begin{equation}\label{eq:physical-rates}
 |\mathcal C_{{\rm CORE},d}|\le5\Lambda^{499/500},\qquad
 |\mathcal C_{{\rm LONG},d}|\le7\Lambda^{19/20}.
\end{equation}
\end{proposition}
\begin{proof}
Use the exact decompositions \eqref{eq:core-five} and
\eqref{eq:long-seven} and apply the twelve physical estimates.
For $\Lambda\ge1$, each CORE exponent is at most $499/500$,
and each LONG exponent is at most $19/20$.
The triangle inequality sums five or seven contributions, so it
costs a factor $5$ or $7$, \emph{not a sum of their exponents}.
The dominant CORE region is $\mathrm{Band}$; the dominant LONG
region is $\mathrm{EdgeL}$. The precise source estimates are
\texttt{core\_pair\_le} and \texttt{long\_pair\_le}.
\end{proof}

The positive exponents in \eqref{eq:physical-rates} are paid for
by explicit geometric inequalities. For the CORE band the
pre-constant exponent is
\begin{equation}\label{eq:band-accounting}
 \underbrace{\frac{199}{200}}_{\text{band width}}
 +\underbrace{\frac1{4000}+\frac1{4000}}_{\text{roughness and weighted mass}}
 =\frac{1991}{2000}
 <\frac{499}{500},
\end{equation}
leaving $1/400$ to absorb a fixed coefficient. For LONG's
square-root edge,
\begin{equation}\label{eq:edge-accounting}
 \underbrace{\frac9{10}}_{\text{thin width}}
 +\underbrace{\frac1{4000}+\frac1{40}+\frac1{100}}_{\text{roughness and summation}}
 =\frac{3741}{4000}<\frac{19}{20},
\end{equation}
leaving $59/4000$. These are checks on the actual cell bounds,
not additional exponent charges after the physical estimates.

\subsection{The decoration cost is charged once}

The full CORE and LONG sectors sum the fixed-decoration pair
sums against $V_h(d)\ge0$:
\begin{equation}\label{eq:decoration-aggregate}
 \mathcal C_{\rm sector}=
 \sum_{d\in\mathcal D_{q,h}}V_h(d)
       \mathcal C_{{\rm sector},d}.
\end{equation}
Thus \eqref{eq:physical-rates} gives the triangle-inequality
bound with \emph{one} factor $\sum_dV_h(d)$. The exact
source decoration estimate is
\begin{equation}\label{eq:decoration-mass}
 \sum_{d\in\mathcal D_{q,h}}V_h(d)
 \le e^{10}\left(1+\frac{\log\Lambda}{\log2}\right)^8
 \le256e^{10}\Lambda^{1/1000}
\end{equation}
for all sufficiently large $\Lambda$. Consequently:
\begin{proposition}[Signed CORE/LONG rates]\label{prop:sector-sublog}
For fixed positive squarefree $q$, $0<\delta<1/3$ and
admissible cutoffs, the complete signed source sectors satisfy
\begin{equation}\label{eq:core-long-final-bounds}
 \boxed{\begin{aligned}
 |\mathcal C_{\rm CORE}|&\le1280e^{10}\Lambda^{999/1000},\\
 |\mathcal C_{\rm LONG}|&\le1792e^{10}\Lambda^{951/1000}.
 \end{aligned}}
\end{equation}
These are eventual bounds uniform in the permitted anchor and
cutoff choices.
\end{proposition}
\begin{proof}
Multiply $5\Lambda^{499/500}$ and
$7\Lambda^{19/20}$ from \eqref{eq:physical-rates} by the
single decoration budget \eqref{eq:decoration-mass}. Then
\[
 \frac{499}{500}+\frac1{1000}=\frac{999}{1000},
 \qquad
 \frac{19}{20}+\frac1{1000}=\frac{951}{1000},
\]
and $5\cdot256=1280$, $7\cdot256=1792$.
No decoration exponent is applied again at terminal assembly.
\end{proof}

The strict, very small $1/1000$ margin in CORE is genuine.
For any $\varepsilon>0$, a large enough $\Lambda$ absorbs
$1280e^{10}\Lambda^{999/1000}$ into
$\varepsilon\Lambda$; LONG has margin $49/1000$.
The source chooses the maximum of the cell threshold, the
polylogarithmic decoration threshold and the constant-absorption
threshold, then a sufficiently large integer $h_0$.
These thresholds are demonstrably finite but potentially enormous.
Neither \eqref{eq:core-long-final-bounds} nor the theorem that
follows claims practically useful estimates at small $h$.

\subsection{The signed Central input}

At the source high-cardinality cutoff $J_u(h)$, the complete HIGH
sector has the eventual estimate
\begin{equation}\label{eq:high-bound}
 |\mathcal C_{\rm HIGH}|\le 8e^{75}(4+48q),
\end{equation}
which is $o_q(\Lambda)$ without another decoration charge.
LOW is exactly zero by \eqref{eq:low-zero}. Taking $\delta=1/4$,
the exact partitions and \eqref{eq:core-long-final-bounds} give
$\mathcal C_{\rm fold}=o_q(\Lambda)$.
The independent correction in \eqref{eq:active-folded} obeys
\begin{equation}\label{eq:folded-error-bound}
 |\mathcal C_{\rm active}-\mathcal C_{\rm fold}|
 \le(6q+2)\,\mathcal E(h),\qquad\mathcal E(h)\longrightarrow0,
\end{equation}
where the source gives an explicit nonnegative reciprocal/density
majorant $\mathcal E(h)$. An $\varepsilon/2$ allowance for the folded
sum and another for the correction prove a uniform signed Central
rate. This is a \emph{Central estimate}, not yet the terminal
inequality for $A^{\rm src}_{1,q}-B^{\rm src}_{1,q}$.

\subsection{The independent shifted Suzuki and Vaaler input}

The other two asymmetric regions are handled by a different analytic
supplier. In its original variables, the sharp shifted Suzuki bound
is uniform for $Y\ge2$, $N\ge1$, $0\le\vartheta\le1$, and
$U\le2N$:
\begin{equation}\label{eq:suzuki-supplier}
 \left|\sum_{N<m\le U}e\!\left(\frac{Y}{m+\vartheta}\right)\right|
 \le C\left(N\exp\!\left[-c\frac{(\log N)^3}{(\log Y)^2}\right]
       +\frac{N^2}{Y}\right),\qquad c>0,\ C\ge0.
\end{equation}
The constants are independent of the varying source parameters.
The proof treats low $Y$ by a discrete telescoping/Kusmin--Landau
argument, and high $Y$ using Karatsuba's selected derivative orders
and Vinogradov mean values \cite{Karatsuba1971}.

The pointwise positive Vaaler input is separately established at the
sawtooth jumps \cite{Vaaler1985}. It bounds a sharply clipped sawtooth
sum by its nonzero Fourier modes \emph{plus a positive error
majorant}. Combined with \eqref{eq:suzuki-supplier} and its
quantified cutoff selection, this supplies the shifted-Karatsuba
bound on the two asymmetric source regions. It cannot be inferred
from Central's CORE/LONG rate.

\subsection{Four complete regions and the native terminal}

For clarity, let $\mathcal Q(q,a,h)$ denote the source's
\emph{literal complete signed quarter-middle} at the original
cutoff $Z(h)$. Its source-matched partition has four regions:
\begin{equation}\label{eq:four-region-identity}
 \boxed{\mathcal Q=\mathcal Q_{\rm left}+\mathcal Q_{\rm right}
    +\mathcal Q_{\rm central}+\mathcal Q_{\rm upper}.}
\end{equation}
These are again complete decorated, signed, quarter-parent sums.
The low-product guard is the exact condition
\[
 \mathsf{LP}_{h,H_d}(r,s):\quad
 rs\le H_d\exp(-\Lambda^{1/4}).
\]
The four mutually exclusive masks are therefore
\[
\begin{array}{ll}
 \mathrm{left}:&\mathsf{LP}\ \wedge\ r^2<s,\\
 \mathrm{right}:&\mathsf{LP}\ \wedge\ s^2<r,\\
 \mathrm{central}:&\mathsf{LP}\ \wedge\ s\le r^2\ \wedge\ r\le s^2,\\
 \mathrm{upper}:&\neg\mathsf{LP}.
\end{array}
\]
The equality cases remain in Central. These masks partition the
entire original quarter-parent population; the definition of
\(\mathsf{LP}\) changes with the reduced height of its decoration.
The names \emph{four-region left/right} refer only to this outer
partition. They use the opposite factor orientation from the
\emph{LONG-left/right} convention within Central: LONG-left means
$s^{1+\delta}<r$ (the $r$ factor is longer), whereas the four-region
left means $r^2<s$ (the $s$ factor is longer). Likewise LONG-right
means $r^{1+\delta}<s$, whereas four-region right means $s^2<r$.
These are distinct partitions, not predicates intended to coincide.

The established shifted Suzuki and positive Vaaler inputs imply the
pointwise $o_q(\Lambda)$ estimates for $\mathcal Q_{\rm left}$
and $\mathcal Q_{\rm right}$; the upper source has its separately
proved rate. The signed Central construction of
\cref{sec:signed-source,sec:cells,sec:aggregation} supplies precisely
$\mathcal Q_{\rm central}=o_q(\Lambda)$ via the source-equality
bridge, with its original cutoffs and decorations. Hence, for each
$\varepsilon>0$, selecting a common threshold for the four source
estimates requested at $\varepsilon/4$ yields
\begin{equation}\label{eq:four-region-budget}
 |\mathcal Q|\le
 |\mathcal Q_{\rm left}|+|\mathcal Q_{\rm right}|+
 |\mathcal Q_{\rm central}|+|\mathcal Q_{\rm upper}|
 \le\varepsilon\log h.
\end{equation}
This triangle inequality is the actual meeting point of Suzuki,
Vaaler, Central and the independently controlled upper region.

There is one further source boundary: $\mathcal Q$ is the
quarter-middle, not the entire exterior. The already-proved
far-quarter estimate gives, after requesting the middle estimate
at $\varepsilon/2$,
\begin{equation}\label{eq:far-quarter-budget}
 |\mathcal Q_{\rm exterior}|\le
 |\mathcal Q_{\rm exterior}-\mathcal Q|+|\mathcal Q|
 \le\frac\varepsilon2\log h+\frac\varepsilon2\log h.
\end{equation}
The exact fixed-$q$ source-to-kernel assembly then transports this
complete exterior estimate to the native one. In particular the
quantified implication now proved is
\begin{equation}\label{eq:native-epsilon}
 \boxed{\forall\,\varepsilon>0\ \exists h_0(q,\varepsilon)\ \forall
 \text{ legal }a,b,h\ge h_0:\quad
 |A^{\rm src}_{1,q}(a;h)-B^{\rm src}_{1,q}(a,b;h)|
 \le\varepsilon\log h.}
\end{equation}
The source proof of this implication factors through the exact
four-region identity, the far-quarter extension, and the
fixed-$q$ kernel transport. The uniform threshold precedes all
endpoint choices; splitting the tolerance does not charge another
logarithmic exponent. The formal compositions are
\texttt{paperB\_terminal\_of\_suzukiCore\_and\_central} and
\texttt{paperB\_terminal\_of\_central\_only}, the latter
supplying the independently proved Suzuki input rather than
assuming it. The printed theorem then follows from the exact
Euler identification \eqref{eq:euler-source-bridge}.

\section{Proof of the transfer theorem and scope}\label{sec:proof-main}

\begin{proof}[Proof of \cref{thm:fixed-q}]
Fix positive squarefree $q$ and $\varepsilon>0$. Choose the
height threshold in \eqref{eq:native-epsilon}. For every
compatible pair of reduced endpoints and every sufficiently
large eligible height, \cref{prop:euler-bridge} gives
\[
\begin{aligned}
 |A_{1,q}(a,b;h)-B_{1,q}(a,b;h)|
 &=|A^{\mathrm{src}}_{1,q}(a;h)
      -B^{\mathrm{src}}_{1,q}(a,b;h)|\\
 &\le\varepsilon\log h.
\end{aligned}
\]
The threshold was chosen before the endpoints; hence it is
uniform in their compatible residue classes. This is the claimed
pointwise $o_q(\log h)$ estimate. The argument is arithmetic
in the singular-series coefficients and invokes no conjecture
about actual consecutive primes.
\end{proof}

\subsection{What has, and has not, been proved}

The theorem is a fixed-conductor, level-one result. The proof
does not control a growing modulus $q$, all ranks of a connected
expansion, a resummation of an infinite rank hierarchy, or the
actual frequencies of consecutive-prime residue pairs. The
last question belongs to the conditional and conjectural
framework of \cite{LemkeOliverSound2016}.

The Boolean contraction identity is elementary but clarifies
which quantity is transferred. The analytic work is the exact
signed-source decomposition combined with estimates tailored
to all its masks; the literal Euler bridges ensure that the
controlled object is the one printed in the theorem, not an
informally analogous kernel. These contributions should be distinguished from the classical
foundations: M\"obius inversion and CRT, Vaaler's pointwise
Fourier majorants \cite{Vaaler1985}, Karatsuba--Vinogradov
exponential-sum methods \cite{Karatsuba1971}, Vaughan's prime-sum
decomposition \cite{Vaughan1980}, and Weil--Estermann
Kloosterman bounds \cite{Weil1948,Estermann1961}.
The period-restricted context has substantial antecedents in
consecutive-prime studies \cite{GoldstonLedoan2011,LemkeOliverSound2016}
and constrained singular-series sums \cite{Kuperberg2025}.
The novel statement claimed here is the particular fixed-$q$
pointwise coefficient transfer and its exact signed-source proof,
not priority over those general techniques.

\subsection{Formal verification and reproducibility}

The accompanying Lean development supplies exact definitions for the
printed Euler products and the signed physical source, and proves the
quantified theorem under its coprimality, compatibility, parity and
positive-height hypotheses. Its two central branches are the literal
Euler identification and the source-native analytic rate. The sealed R0003 source and reproduction instructions are maintained
separately. On 2 November 2026, Alex Breakell independently reproduced
the pinned-source Lean build on Ubuntu/WSL2 (9,599 jobs, status 0) and
the printed-terminal axiom check (status 0; only standard Lean axioms).
The author thanks him as independent verification reviewer; this is
not independent handwritten mathematical or journal peer review. The analytic appendices in this article develop four
worked mechanism chains with explicit intermediate inequalities,
provides the complete physical-cell contract register, and
tracks the exact costs in the eventual exponent accounting.
It must not be confused with a second independently handwritten
proof of all deep supplier lemmas: weighted rough-set estimates,
source-specific completion bounds and reciprocal-source
estimates remain identified, precisely scoped mathematical inputs
from the sealed formal development.

\paragraph{AI-assisted research and verification.}
Generative AI systems contributed extensively to mathematical research,
Lean proof development, manuscript drafting and internal critical review
under human direction. These AI-assisted referee exercises do not constitute
independent human journal peer review. The named human author is responsible
for the final scholarly content, attributions and public submission.
The status of the source build and the limits of the formal evidence are
reported above and in the separately maintained release record.

\medskip
\appendix
\phantomsection
\pdfbookmark[1]{Analytic proof appendices}{technical-appendices}
\section*{Analytic proof supplement: four mechanisms}
These appendices expand the source applications singled out in the main text,
identify their deep imported estimates precisely, and preserve the exact
hypotheses under which those estimates are used. They integrate the separately
reviewed analytic supplement into one mathematical publication. Source-locked
formal declarations and a declaration-level obligation ledger permit the
individual imports to be inspected; this editorial integration does not itself
constitute another Lean replay.
\section{What is being proved, and at what evidence level}\label{sec:contract}

Fix an integer $q>0$ that is squarefree. The anchor $a$ is coprime
to $q$, the height $h>0$ satisfies the required parity when $2\nmid q$,
and the second endpoint class is compatible with $a+h$. Put
\begin{equation}\label{eq:basic-parameters}
 \Lam=\log h,\qquad \sigma_0=\frac23+\frac1{100},\qquad
 \epsp=\Lam^{-1/200},\qquad \eta=\Lam^{-1/10},\qquad
 0<\delta<\frac13.
\end{equation}
The numbers $\epsp$ and $\eta$ are \emph{prescribed internal functions
of $h$}; neither is the arbitrary accuracy parameter in the theorem's
$\varepsilon\log h$ conclusion. The final proof takes $\delta=1/4$.
Everything below is in the eventual range of sufficiently large $\Lam$.

Let $d$ be a legal decoration by primes dividing $h$, and put
\[
 D=\prod_{p\in d}p,\qquad H=h/D,\qquad
 V_h(d)=\prod_{p\in d}\frac1{p-2}.
\]
The original positive and high-prime restrictions are the ones in
Paper B, not a freely chosen replacement: $Z(h)=2^{20(\lfloor
\Lam^{\sigma_0}\rfloor+1)}$, and
\begin{equation}\label{eq:physical-universe}
 T=T_{q,h,d}=\{p\le H^2: p\text{ prime},\ p\ge5,
       \ p\nmid qH,\ p\notin d\},\qquad
 \omegaR(B)=\prod_{p\in B}(p-3)^{-1}.
\end{equation}
The \emph{ordered} pair $(B,C)$ consists of nonempty disjoint subsets
of $T$, $r=\prod_{p\in B}p$, $s=\prod_{p\in C}p$,
subject to the quarter-parent bounds
$Z(h)<rs$, $4rs\le H^2$, the cardinality window
$J_{\rm low}(h)\le |B|+|C|<J_{\rm high}(h)$, and the original Central gate
\begin{equation}\label{eq:central-gate}
 rs\le H\exp(-\Lam^{1/4}),\qquad s\le r^2,\qquad r\le s^2.
\end{equation}
Let $\F_{q,a,D}(H,r,s)$ denote the \emph{same ordered signed
folded row} as in the main paper, including its reciprocal sawtooth
and the ordered difference of its two CRT potentials. The sealed
source establishes the uniform envelope
\begin{equation}\label{eq:row-envelope}
 |\F_{q,a,D}(H,r,s)|\le F_q:=4+48q.
\end{equation}
All results in this supplement involve the complete masked sum
\begin{equation}\label{eq:masked-sum}
 \mathcal P_d(\mathsf P)=
 \sum_{\substack{B,C\subseteq T\ \,B\cap C=\varnothing\\B,C\ne\varnothing}}
 \one_{\MR(B,C)\wedge\mathsf P(B,C)}\,
 \omegaR(B)\omegaR(C)\F_{q,a,D}(H,r,s),
\end{equation}
where $\MR$ retains every preceding source gate and $\mathsf P$
is the cell-specific predicate. In particular, a bound for
$\mathcal P_d(\mathsf P)$ does \emph{not} assert a bound of the same
size for every individual parent. For the two sectors, the source
exactly proves
\begin{align}
\mathcal C_{\csmall,d}&=W_d+C_d+R_d+\mathrm{ResA}_d+\mathrm{Band}_d,\notag\\
\mathcal C_{\clong,d}&=WL_d+CL_d+RL_d+\mathrm{ManyL}_d\notag\\
&\quad+\mathrm{EdgeL}_d+\mathrm{ResL}_d+WRL_d.
\label{eq:sector-identity}
\end{align}
Every summand on the right is an actual instance of
\eqref{eq:masked-sum}; the identities use predicate complements
\emph{before} any triangle inequality.

\begin{remark}[What is proved on these pages]\label{rem:standing}
The logical implications from the masks to geometric intervals,
the exact logarithmic exponent calculations, the norm inequalities
following from specified auxiliary bounds, and the finite
centering algebra are derived below. Deep supplier estimates such
as the weighted rough-set lemma, harmonic interval-mass lemma,
physical reciprocal source estimate and the coupled Weil--Cauchy
bound are \emph{imported mathematical lemmas} whose complete proofs
reside in the sealed Lean modules. They are stated with explicit
hypotheses and source positions so that a conventional referee can
attack the right interface. Reproducing a theorem's numerical
exponent here is not a second independent proof of that supplier.
\end{remark}

\section{The CORE Band: deriving the one-thousandth saving}\label{sec:band}

In the CORE mask, write $\mathsf W(B,C)$ for the B-factor
prime-packet test, and use the source's packet-width parameter
$\epsp=\Lam^{-1/200}$. After excluding the earlier $W$ and $C$
cells, the remaining decision distinguishes the $R/\mathrm{ResA}$
ranges from the Band range. The Band predicate, in addition to the
complete source selection in \eqref{eq:masked-sum}, is
\begin{equation}\label{eq:band-mask}
 \neg\mathsf W(B,C),\qquad
 \neg(s^{1+\epsp}\le r),\qquad
 \neg(r^{1+3\epsp}\le s).
\end{equation}
These are \emph{strict} complements of the preceding weak
inequalities. No positivity or endpoint condition may be dropped
in passing to the next statement.

\begin{lemma}[The band is a short logarithmic interval]
\label{lem:band-interval}
For a Band-selected pair, $r,s>1$ and
\begin{equation}\label{eq:band-window}
 r^{1/(1+\epsp)}<s<r^{1+3\epsp}.
\end{equation}
For $A=r^{1/(1+\epsp)}$, $B_+=r^{1+3\epsp}$,
$a_0=\lfloor A\rfloor$, $b_0=\lfloor B_+\rfloor$,
the integer $s$ belongs to $(a_0,b_0]$ and, in the eventual
range where $A\ge2$,
\begin{equation}\label{eq:band-ratio}
 \frac{b_0}{a_0}\le2r^{4\epsp},\qquad
 \log\frac{b_0}{a_0}\le1+4\Lam^{199/200}.
\end{equation}
\end{lemma}
\begin{proof}
The first negation in \eqref{eq:band-mask} involving $s^{1+\epsp}$
gives $r<s^{1+\epsp}$; raising both positive sides to
$1/(1+\epsp)$ gives the first inequality in
\eqref{eq:band-window}. The other negation gives the second.
For $A\ge2$, $\lfloor A\rfloor\ge A/2$ and
$\lfloor B_+\rfloor\le B_+$. Thus
\[
 \frac{b_0}{a_0}\le 2\frac{B_+}{A}
  =2r^{1+3\epsp-1/(1+\epsp)}\le2r^{4\epsp},
\]
since $1-1/(1+u)=u/(1+u)\le u$ for $u>0$.
The source gates imply $r\le h$, so $\log r\le\Lam$.
Since $\epsp\Lam=\Lam^{199/200}$ and
$\log2<1$, taking logarithms yields
\eqref{eq:band-ratio}. The floor endpoints are retained:
no continuous approximation has replaced the discrete interval.
\end{proof}

The relevant weighted counting supplier is much more precise
than a bound by the number of integer points. For a finite set
$T$ of primes at least five, let
\(M_T(a,b)=\sum_{C\subseteq T:\,a<\prod C\le b}\omegaR(C)\).
The source's harmonic interval-mass lemma gives, for $1\le a\le b$,
\begin{equation}\label{eq:interval-mass-supplier}
 M_T(a,b)\le e^{37.5}
 \bigl(1+\log_2\log_2 b\bigr)^{30}
 \bigl(1+\log(b/a)\bigr),
\end{equation}
where the double binary logarithm is the integer \texttt{Nat.log}
operation, not an unqualified real logarithm. The proof uses
injectivity of the prime-subset product and the harmonic interval
bound after a pointwise omega majorant. For the present interval,
$\log b_0\le5\Lam$ eventually. The source proves that the fixed
polylogarithmic prefactor in \eqref{eq:interval-mass-supplier}
is then at most $\Lam^{1/4000}$. Consequently,
\begin{equation}\label{eq:band-C-bound}
 \sum_{C:\,r^{1/(1+\epsp)}<s<r^{1+3\epsp}}\omegaR(C)
 \le \Lam^{1/4000}(2+4\Lam^{199/200})
 \le 6\Lam^{1/4000+199/200}.
\end{equation}
The $C$-sum here is over subsets of the actual $T$, and
\eqref{eq:band-C-bound} is an \emph{upper bound} after dropping
other source gates, not an equality of masks.


\begin{lemma}[The missing $W$ packet pays for roughness]
\label{lem:band-rough-transport}
Write $Y_p=\lfloor s^{u_0(\delta)}\rfloor$ for the $W$-packet upper endpoint,
where $u_0(\delta)>0$, and set $u=u_0(\delta)/2$.
On the Band mask, every $p\in B$ satisfies
\begin{equation}\label{eq:band-rough-bridge}
 p<X_p(h)\quad\text{or}\quad p>r^u.
\end{equation}
Moreover the Band-specific product inequalities imply, eventually,
\begin{equation}\label{eq:band-five-bridge}
 \log r>\frac{20\log2}{2+3\epsp}\Lam^{\sigma_0}
    \ge5\Lam^{\sigma_0},\qquad
 r\ge\left\lceil\exp(5\Lam^{\sigma_0})\right\rceil.
\end{equation}
\end{lemma}
\begin{proof}
If the $W$ packet is absent, no prime in $B$ lies in the integer
interval $[X_p(h),Y_p]$. Thus a prime not below $X_p(h)$ satisfies
$p>\lfloor s^{u_0}\rfloor$; since $p$ is an integer, this implies
$p>s^{u_0}$. The Band inequality $r<s^{1+\epsp}$ and
$\epsp\le1$ give $r<s^2$, hence
$s^{u_0}\ge r^{u_0/2}=r^u$. This proves
\eqref{eq:band-rough-bridge} without replacing a missing packet by
an assumed source prime.

For the height threshold, the strict quarter-parent bound and the
upper Band inequality give
\[
 Z(h)<rs<r^{2+3\epsp},\qquad
 \log Z(h)=20\bigl(\lfloor\Lam^{\sigma_0}\rfloor+1\bigr)\log2
    >20(\log2)\Lam^{\sigma_0}.
\]
Division by $2+3\epsp>0$ proves the first inequality of
\eqref{eq:band-five-bridge}. The coefficient tends to
$10\log2>5$ as $\epsp\to0$, so it is at least $5$ eventually.
The ceiling follows because $r$ is an integer. This stronger
$5$-threshold uses the \emph{refined Band interval}; the generic CORE
gate $rs\le r^3$ alone does not yield it.
\end{proof}

In applying \eqref{eq:interval-mass-supplier}, the integer interval
has $1\le a_0\le b_0$ in the eventual source range. It counts
\emph{nonnegative} subset-product weights over the actual finite prime
set, including the empty subset if its product lies in the interval;
the source masks only restrict this population further. The rough
supplier \eqref{eq:band-B-bound} is applied to products $r$ in the
actual high-support window
$\lceil e^{5\Lam^{\sigma_0}}\rceil\le r\le h$,
with $u=u_0(\delta)/2$ \emph{fixed} after $\delta$ is chosen.
The matching source contracts are
\path{QuillonResidue.interval_mass},
\path{QuillonResPhysical.rough_ev}, and
\path{QuillonResPhysical.Band_physical}.

The missing W packet also makes the \emph{$B$-factor} arithmetically
sparse. Put $u=u_0(\delta)/2>0$ as in the source rough-set
lemma and $R_0=\lceil e^{5\Lam^{\sigma_0}}\rceil$.
On the Band-selected population the source proves
\[
 R_0\le r\le h,\qquad
 \forall p\in B:\ p<X_p(h)\ \text{or}\ r^u<p.
\]
The weighted roughness lemma (with its explicit exponential
eligibility condition) gives
\begin{equation}\label{eq:band-B-bound}
 \sum_{B:\,\text{these roughness tests}}\omegaR(B)
 \le \Lam^{1/4000}(1+u^{-5}).
\end{equation}
Here $u>0$ is fixed when $\delta$ is fixed, while $\epsp$
is a shrinking function of $h$; exchanging the two would
change the argument.

\begin{proposition}[Band estimate from its two supplier lemmas]
\label{prop:band-chain}
Assume \eqref{eq:row-envelope}, the source-valid mask implications
above, and the interval and roughness bounds
\eqref{eq:band-C-bound}--\eqref{eq:band-B-bound}.
Then, for sufficiently large $\Lam$,
\begin{equation}\label{eq:band-chain}
 |\mathrm{Band}_d|
 \le K_{q,\delta}\Lam^{1991/2000}
 \le\Lam^{499/500},\qquad
 K_{q,\delta}=6(4+48q)(1+u^{-5}).
\end{equation}
\end{proposition}
\begin{proof}
Apply the absolute triangle bound to the double sum
\eqref{eq:masked-sum}, then enlarge the $C$-mask to the interval
\eqref{eq:band-window} and the $B$-mask to the roughness
set. The weights are nonnegative. Factoring the two independent
upper bounds yields
\[
 |\mathrm{Band}_d|
 \le F_q\Bigl[\Lam^{1/4000}(1+u^{-5})\Bigr]
       \Bigl[6\Lam^{1/4000+199/200}\Bigr]
 =K_{q,\delta}\Lam^{1991/2000},
\]
because $2/4000+199/200=1991/2000$. Finally
\[
 \frac{499}{500}-\frac{1991}{2000}=\frac1{400}>0.
\]
For fixed $q,\delta$, the constant $K_{q,\delta}$ is absorbed
once $\Lam^{1/400}\ge K_{q,\delta}$. This is precisely
the source's final \texttt{const\_le\_rpow} step, not an
independent loss introduced by summing decorations.
\end{proof}

This calculation identifies the bottleneck: the signed CORE
sum is eventually dominated by a \emph{positive} band-mass
bound, even though its other source cells have strong negative
powers. The margin available after full sector aggregation is
smaller still, as shown in \cref{sec:accounting}.

\section{The LONG square-root edge and the Weil region}\label{sec:edge}

The LONG predicate is
\[
 \mathsf{Long}_{\delta}(r,s)\iff
 s^{1+\delta}<r\ \text{or}\ r^{1+\delta}<s.
\]
The first orientation is called LONG-left; the second LONG-right.
These are \emph{not} the left/right regions of the earlier
four-region quarter-middle decomposition. Here we work only on
LONG-right and retain the original common mask $\MR_L$.

Write
\begin{equation}\label{eq:jr-def}
 \JR(h,T,B)=\{p\in T\setminus B:
 X_p(h)<p\le 2^{K_v(h,r)}X_p(h)\},
\qquad K_0(h)=\lceil100\log\Lam\rceil.
\end{equation}
On LONG-right the source tests in order whether $C\cap\JR$
is nonempty, whether $|C|\ge K_0$, and then whether
\begin{equation}\label{eq:edge-mask}
 r<s^{1/2+3\eta}.
\end{equation}
If the first test is true the pair belongs to $RL$; if the
second is true it belongs to $\mathrm{ManyL}$. Only the
remaining pairs with \eqref{eq:edge-mask} enter $\mathrm{EdgeL}$.
Off the edge, the presence or absence of a B-factor packet
\[
 \mathcal W_L(B,C)=
 B\cap\{p\in T\setminus C:X_p(h)\le p\le\lfloor s^\eta\rfloor\}
\]
selects $WRL$ or $\mathrm{ResL}$, respectively. This ordered
routing is exact, not an overlap of estimates.

\subsection{EdgeL: a geometric estimate rather than cancellation}

\begin{lemma}[Geometric edge strip]\label{lem:edge-strip}
On the full EdgeL mask, the Central gate $s\le r^2$ and
\eqref{eq:edge-mask} imply
\begin{equation}\label{appendix:eq:edge-strip}
 s^{1/2}\le r<s^{1/2+3\eta}.
\end{equation}
Fixing $C$ (and hence $s$) puts the possible $B$-products in
\[
 r\in\bigl(\lfloor s^{1/2}\rfloor-1,\,
                   \lfloor s^{1/2+3\eta}\rfloor\bigr],
\]
with logarithmic width at most $1+3\Lam^{9/10}$
up to the explicit endpoint-floor adjustment.
\end{lemma}
\begin{proof}
Since $r,s\ge1$, the gate $s\le r^2$ implies
$s^{1/2}\le r$. The opposite bound is exactly the EdgeL
test. Put $u=s^{1/2}$ and $v=s^{1/2+3\eta}=u s^{3\eta}$.
For sufficiently large $h$ in the selected source $u\ge8$;
thus $a_0=\lfloor u\rfloor-1\ge u/2>0$ and
$b_0=\lfloor v\rfloor\le v$. Consequently
\[
 b_0/a_0\le2s^{3\eta},\qquad
 \log(b_0/a_0)\le\log2+3\eta\log s
 \le1+3\Lam^{9/10},
\]
using $s\le h$. If $r\ge u$ is an integer then
$r>\lfloor u\rfloor-1$; and $r<v$ implies
$r\le\lfloor v\rfloor$. The indicated interval therefore
contains every EdgeL B-product.
\end{proof}


\begin{lemma}[An absent reciprocal packet gives $B$-independent rough $C$]
\label{lem:edge-rough-transport}
On the ordered EdgeL mask, suppose $C\subseteq T\setminus B$,
$C\cap\JR(h,T,B)=\varnothing$, $r\ge1$ and $s\le r^2$.
For every $p\in C$ the following dichotomy holds:
\begin{equation}\label{eq:edge-rough-bridge}
 p<X_p(h)+1\qquad\text{or}\qquad p>s^{\epsp/2}.
\end{equation}
It is independent of $B$ once $s$ and $h$ are fixed.
\end{lemma}
\begin{proof}
For a prime $p\in C$ with $p\le X_p(h)$, the first alternative
is immediate. Otherwise $p>X_p(h)$; since $p\notin\JR(h,T,B)$,
its integer packet upper endpoint is exceeded:
\[
 p>2^{K_v(h,r)}X_p(h)\ge r^{\epsp}
       \ge s^{\epsp/2}.
\]
Here $X_p(h)\ge1$, the source's definition of $K_v$ guarantees
$2^{K_v(h,r)}\ge r^{\epsp}$, and the last inequality follows from
$s\le r^2$ with $\epsp>0$. Consequently the source's
\texttt{rough\_ev4} is applied to a genuine $C$-only roughness
population, not to a family still depending on the varying $B$.
This is the mathematical content of
\texttt{QuillonLRes.C\_rough\_of}.
\end{proof}

The integer lower endpoint in \cref{lem:edge-strip} is likewise
paid by the actual selected source: on LONG-right $r<s$ and
$Z(h)<rs<s^2$, so $s>\sqrt{Z(h)}$ and eventually
$\sqrt s\ge8$. Thus
$\lfloor\sqrt s\rfloor-1\ge1$ and the harmonic interval-mass
lemma applies to its prescribed positive integer interval. The
roughness sum for $C$ then has source lower support
$\lceil e^{4\Lam^{\sigma_0}}\rceil\le s\le h$,
not the stronger CORE $5$-threshold.

The absent reciprocal packet makes $C$ rough in the precise
sense required by the LONG source. The source's $\texttt{rough\_ev4}$
then yields, at $R_0=\lceil e^{4\Lam^{\sigma_0}}\rceil$,
\begin{equation}\label{eq:edge-C-bound}
 \sum_{C:\,\mathrm{rough}}\omegaR(C)
 \le\Lam^{1/4000}\bigl(1+(\epsp/2)^{-5}\bigr)
 =\Lam^{1/4000}(1+32\Lam^{1/40})
 \le33\Lam^{1/4000+1/40}.
\end{equation}
The equality $2^5\Lam^{5/200}=32\Lam^{1/40}$ is a
real part of the quantitative cost, not a negligible constant.

Fix one such $C$. Applying the same harmonic interval-mass
supplier \eqref{eq:interval-mass-supplier}, now to the
$B$-products in \cref{lem:edge-strip}, and using the
source's bound $e^{37.5}(1+\log_2\log_2 b_0)^{30}\le
\Lam^{1/100}$, yields
\begin{equation}\label{eq:edge-B-bound}
 \sum_{B:\,s^{1/2}\le r<s^{1/2+3\eta}}\omegaR(B)
 \le\Lam^{1/100}(2+3\Lam^{9/10})
 \le5\Lam^{1/100+9/10}.
\end{equation}
Again the enlargement forgets restrictions, never counts a
term with a negative weight, and is valid because $\omegaR(B)\ge0$.

\begin{proposition}[EdgeL estimate from geometric mass]
\label{prop:edge-chain}
Assume the source-valid roughness and interval supplier bounds
\eqref{eq:edge-C-bound}--\eqref{eq:edge-B-bound} and
\eqref{eq:row-envelope}. Then
\begin{equation}\label{eq:edge-chain}
 |\mathrm{EdgeL}_d|\le165(4+48q)
 \Lam^{3741/4000}\le\Lam^{19/20}
\end{equation}
eventually, with no Kloosterman input.
\end{proposition}
\begin{proof}
The proof reverses the order used for Band: fix $C$ first,
apply the geometric interval bound to the sum over $B$,
then sum over rough $C$. One obtains
\[
 |\mathrm{EdgeL}_d|
 \le F_q[33\Lam^{1/4000+1/40}]
          [5\Lam^{1/100+9/10}]
 =165F_q\Lam^{3741/4000},
\]
where
\[
 \frac1{4000}+\frac1{40}+\frac1{100}+\frac9{10}
 =\frac{3741}{4000},\qquad
 \frac{19}{20}-\frac{3741}{4000}=\frac{59}{4000}.
\]
A fixed-coefficient absorption at exponent $59/4000$
completes the stated physical bound. There is no cancellation
estimate anywhere in this derivation beyond the source's
weighted mass suppliers.
\end{proof}

\subsection{WRL: why the same interval cannot be treated by Weil}

The off-edge test is $r\ge s^{1/2+3\eta}$, while $WRL$ requires
an actual prime $p\in\mathcal W_L(B,C)$.
That prime satisfies $p\le s^\eta$, so extraction leaves
\begin{equation}\label{eq:weil-cofactor}
 \frac rp\ge s^{1/2+2\eta}.
\end{equation}
The $s^{2\eta}$ gain above the square-root length is the room
used by the source's completed bilinear argument; it is
\emph{not} present uniformly on the edge strip
\eqref{appendix:eq:edge-strip}. Merely asserting the existence of
a strong complete-sum estimate cannot replace this range
hypothesis.

For a complete Kloosterman sum with the source coefficient
convention, the genuine supplier is the general-modulus bound
\begin{equation}\label{eq:weil-contract}
 |K_s(u,v)|\le\tau(s)\sqrt{s\,\gcd(u,v,s)}.
\end{equation}
The equality transporting the physical $K_s$ to the general
Kloosterman normalization reverses $u,v$. In the actual
squarefree case $s=\prod_{p\in C}p$, $\tau(s)=2^{|C|}$;
for the frequency-variable modulus the bound retains the
product of primes annihilating the Fourier coefficient.
The conductor-gcd penalty is not set to $1$.

The coupled Cauchy estimate splits a square of a Fourier-row
sum into three contributions: $y=y'$, $y\ne y'$ but
$y\equiv y'\pmod s$, and the genuinely noncongruent pairs.
In the third case the source has a factor of the form
\[
 4^{|C|/2}s^{1/2}G_0^{1/2}
 \bigl(L/s+1+\log s\bigr)
 (1+\log Y)^2
 \prod_{p\in C}(1+p^{-1/2}),
\]
where $G_0$ is the product of the source primes dividing the
fixed phase coefficient. After taking the Cauchy square root,
a factor $G_0^{1/4}$ remains and has to be summed using a
\emph{separate} harmonic gcd-frequency moment. The row
has four source Fourier contributions (two main and two
class terms); all four must be controlled. These are genuine
analytic obligations, not descriptive decoration.


\begin{remark}[The discrete cofactor and the Weil hypothesis are paid]
\label{rem:wrl-discrete-weil}
The real inequality \eqref{eq:weil-cofactor} does not identify the
\emph{integer} range in the subsequent Ramar\'e decomposition.
For $Y_\eta=\lfloor s^\eta\rfloor$ and the source's lower
cofactor cutoff $L_{w,\eta}$, the actual integer supplier gives,
provided $0<\eta\le1/8$ and $s^\eta\ge8$,
\begin{equation}\label{eq:wrl-integer-cofactor}
 1\le\left\lfloor\frac{L_{w,\eta}}{Y_\eta}\right\rfloor,
 \qquad
 \frac12s^{1/2+2\eta}
 \le\left\lfloor\frac{L_{w,\eta}}{Y_\eta}\right\rfloor.
\end{equation}
The explicit factor $1/2$ protects the floor; it is not an
unjustified replacement of the integer cofactor by $r/p$.
The source's named \texttt{WeilHyp} is not a terminal axiom.
It is supplied by \path{Orienne.CSC3.quillon_weilHyp} via
\path{norm_quillon_K_weil}, whose exact identity is
$K_s(u,v)=\mathrm{kSum}_s(v,u)$, reversing the two coefficients.
The general-modulus input \eqref{eq:weil-contract} retains both
$\tau(s)$ and the full conductor $\gcd(u,v,s)$.
\end{remark}

\begin{lemma}[The two conductor moments retained in WRL]
\label{lem:wrl-moments}
Let $P$ be a finite set of primes and let $Y,J\ge1$.
The source's nonnegative harmonic-gcd estimates are
\begin{align}
&\sum_{1\le y\le Y}\ \sum_{\substack{1\le y'\le Y\\y'\ne y}}
 \frac1{yy'}
 \left(\prod_{\substack{p\in P\\p\mid(y-y')}}p\right)^{1/2}
 \le 2(1+\log Y)^2\prod_{p\in P}(1+p^{-1/2}),
 \label{eq:wrl-pair-moment}\\
&\sum_{1\le j\le J}\frac1j
 \left(\prod_{\substack{p\in P\\p\mid j}}p\right)^{1/4}
 \le(1+\log J)\prod_{p\in P}(1+p^{-3/4}).
 \label{eq:wrl-mode-moment}
\end{align}
Both bounds are supplied by the proved source lemmas
\texttt{gcd\_moment\_half} (pair differences) and
\texttt{gcd\_mode\_sum} (mode weights at exponent $1/4$);
we do not re-prove those suppliers here.
\end{lemma}

\noindent
For clarity, the source's completed proof expends these bounds in
\emph{two separate places}. After decomposing the square into equal,
distinct-congruent and noncongruent pairs, the off-diagonal
$\gcd$ majorant contains a square-root product of primes dividing
the pair difference. Equation \eqref{eq:wrl-pair-moment} pays that
weighted pair difference. Taking the Cauchy square root retains
$G_0^{1/4}$, and \eqref{eq:wrl-mode-moment} pays the surviving
\emph{mode}-frequency weight rather than deleting it.
Positive and negative main modes and both class modes are bounded
separately with their actual coefficients. The two moments are
nonnegative upper bounds; neither is asserted to equal the source
signed sum. The physical \texttt{W\_final\_LR} leaf additionally
proves the Fourier-window and coefficient conditions before yielding
\eqref{eq:wrl-fixed-C}; the displayed moments alone are not a
new proof of that deep leaf.

The fixed-$C$ physical theorem applies under the proved
Weil hypothesis and the required range, mode, and coefficient
conditions. It gives
\begin{equation}\label{eq:wrl-fixed-C}
 \left|\sum_{\substack{B\subseteq T\setminus C\\\mathsf{WRL}(B,C)}}
 \omegaR(B)\F_{q,a,D}(H,r,s)\right|
 \le4\Lam^{-95}.
\end{equation}
This \emph{deep analytic leaf} is supplied by
\texttt{GoldC.QuillonLWR.W\_final\_LR}; the preceding
paragraphs expose its key mechanism but do not rederive
its full completion argument.
The remaining $C$-factor is nonnegative and the source's
outer mass bound is
\begin{equation}\label{eq:outer-mass}
 \sum_{C\subseteq T}\omegaR(C)
 \le e^{10}5^{10}\Lam^{10}.
\end{equation}
Thus $|WRL_d|\le4e^{10}5^{10}\Lam^{-85}$.
For fixed constants and sufficiently large $\Lam$,
$4e^{10}5^{10}\le\Lam$ and $\Lam^{-84}\le\Lam^{-80}$,
so $|WRL_d|\le\Lam^{-80}$. The packet-absent complement
belongs to $\mathrm{ResL}$ and is paid by rough mass, not
by applying \eqref{eq:wrl-fixed-C} without its packet.

\section{R and RL: the same reciprocal sum in two geometries}
\label{sec:reciprocal}

The moving prime packet in the $C$-factor is the actual set
\begin{equation}\label{eq:Rpacket}
 \JR(h,T,B)=\{p\in T\setminus B:\quad
 X_p(h)<p\le2^{K_v(h,r)}X_p(h)\},
 \quad K_v(h,r)=\lfloor\log_2\lceil r^{\epsp}\rceil\rfloor+1.
\end{equation}
The CORE $R$ mask requires
\[
 \neg\mathsf W(B,C),\quad
 \neg(s^{1+\epsp}\le r),\quad
 r^{1+3\epsp}\le s,\quad
 C\cap\JR(h,T,B)\ne\varnothing,
\]
with the CORE source gate. LONG-right's $RL$ mask instead
requires $r^{1+\delta}<s$ and the same nonempty reciprocal
packet, together with the source's LONG mask. Eventually
$3\epsp\le\delta$, so
\begin{equation}\label{eq:RL-geometry}
 r^{1+3\epsp}\le r^{1+\delta}<s.
\end{equation}
The geometry is enough for a common reciprocal strategy,
but \emph{not} permission to reuse every CORE numerical
hypothesis in LONG.

\subsection{How the shifted Suzuki input enters the physical row}

For $Y\ge2$, $N\ge1$, $0\le\vartheta\le1$ and $U\le2N$,
the proved shifted reciprocal input is of the form
\begin{equation}\label{appendix:eq:suzuki-supplier}
 \left|\sum_{N<m\le U}e\left(\frac{Y}{m+\vartheta}\right)\right|
 \le C\left[N\exp\left(-c\frac{(\log N)^3}{(\log Y)^2}\right)
            +\frac{N^2}{Y}\right],\qquad e(x)=e^{2\pi ix},
\end{equation}
for fixed $c>0,C\ge0$ independent of the displayed
variables. The small-$N$ and low-frequency cases are
covered by the source's separate branch; no high-frequency
hypothesis is silently imposed on all parameters.

The physical residue class uses an integer modulus $Q=6qr$
(when $r=\prod B$). For a residue $0\le\rho<Q$,
writing $m=Qn+\rho$ gives the exact phase identity
\begin{equation}\label{eq:ap-phase}
 \frac{Y}{Qn+\rho}=
 \frac{Y/Q}{n+\rho/Q},\qquad
 \vartheta=\rho/Q\in[0,1).
\end{equation}
Only after the transformed $Y/Q$ and clipped interval satisfy
the range hypotheses of \eqref{appendix:eq:suzuki-supplier} may the
latter be invoked. A negative numerator is handled through
$e(-x)=\overline{e(x)}$, which preserves absolute value.
The Fourier coefficient attached to the source mode is
not constant. If $S_t=\sum_{n=N+1}^t z_n$, Abel summation
furnishes the exact identity
\[
 \sum_{n=N+1}^{U}c_nz_n
 =c_US_U+\sum_{t=N+1}^{U-1}(c_t-c_{t+1})S_t
\]
and therefore
\begin{equation}\label{eq:abel-cost}
 \left|\sum_{n=N+1}^{U}c_nz_n\right|
 \le \left(|c_U|+\sum_{t=N+1}^{U-1}|c_t-c_{t+1}|\right)
     \max_{N<t\le U}|S_t|.
\end{equation}
Thus the physical Suzuki application spends both a
progression translation and a \emph{variation norm} of the
actual coefficients. Using \eqref{appendix:eq:suzuki-supplier} as if
the source coefficients were all $1$ would be invalid.


\begin{lemma}[The clipped, weighted $R/RL$ Suzuki application]
\label{lem:clipped-suzuki}
Let $Q\ge1$, $0\le\rho<Q$, $V_1\ge1$ and
$V_1<v,v'\le2V_1$ with $v\ne v'$. Fix the source's outer
interval $(U_0,U_1]$, its product clip $(S_0,S_1]$, and
$j,L\ge1$. Assume the \emph{actual source inequalities}
\begin{equation}\label{eq:suzuki-source-window}
 4Q\le U_0,\qquad U_1+2Q\le2U_0,
 \qquad |Y|\ge8QV_1^2.
\end{equation}
Set, using integer divisions,
\begin{align}
 a_*&=\max\!\left\{U_0,\left\lfloor S_0/v\right\rfloor,
                    \left\lfloor S_0/v'\right\rfloor\right\},
 & b_*&=\min\!\left\{U_1,\left\lfloor S_1/v\right\rfloor,
                    \left\lfloor S_1/v'\right\rfloor\right\},
 \label{eq:suzuki-actual-clip}\\
 N&=(a_*-\rho)\mathbin{\mathrm{div}}Q,
 &M&=(b_*-\rho)\mathbin{\mathrm{div}}Q,\quad
 \vartheta=\rho/Q,
 \qquad Y'=Y/v-Y/v'.
 \label{eq:suzuki-transformed}
\end{align}
Here $\mathrm{div}$ is natural-number quotient, as in the
source. On every nonempty clipped residue interval,
\begin{equation}\label{eq:suzuki-eligibility}
 4Q\le a_*,\qquad b_*+2Q\le2a_*,\qquad
 1\le N<M\le2N,\qquad
 |Y'|\ge\frac{|Y|}{4V_1^2}\ge2Q.
\end{equation}
For $t\in(N,M]$ let
\[
 A_t=\mathrm{ampC}\bigl(j,vL(\rho+Qt)\bigr)\,
     \mathrm{ampC}\bigl(j,v'L(\rho+Qt)\bigr).
\]
The actual cotangent amplitudes satisfy
$0\le A_t\le1$ and are nondecreasing in $t$, so
\begin{equation}\label{eq:suzuki-amplitude-variation}
 \mathcal V(A):=|A_M|+\sum_{t=N+1}^{M-1}|A_{t+1}-A_t|
     =2A_M-A_{N+1}\le2.
\end{equation}
For the constants $c>0,C\ge0$ of the source's unconditional
shifted Suzuki theorem,
\begin{equation}\label{eq:suzuki-clipped-conclusion}
 \left|\sum_{t=N+1}^M A_t
 e\!\left(\frac{Y'/Q}{t+\vartheta}\right)\right|
 \le2C\left[N\exp\left(
 -c\frac{(\log N)^3}{(\log(|Y'|/Q))^2}\right)
       +\frac{N^2}{|Y'|/Q}\right].
\end{equation}
The empty interval contributes zero instead.
\end{lemma}
\begin{proof}
The two product constraints $S_0<vu,v'u\le S_1$ and the
outer clip $U_0<u\le U_1$ give exactly $(a_*,b_*]$.
Because $a_*\ge U_0$ and $b_*\le U_1$, the first two
inequalities in \eqref{eq:suzuki-eligibility} follow from
\eqref{eq:suzuki-source-window}. For a nonempty residue interval,
write $u=\rho+Qt$; natural division gives the exact endpoints
$N,M$. The inequalities $4Q\le a_*$ and
$b_*+2Q\le2a_*$ imply $N\ge1$ and $M\le2N$
with the floor remainders retained. The difference phase satisfies
\[
 \frac{Y}{vu}-\frac{Y}{v'u}=\frac{Y'}{u},\qquad
 |Y'|=\frac{|Y|\,|v-v'|}{vv'}
       \ge\frac{|Y|}{4V_1^2};
\]
we used $|v-v'|\ge1$ for distinct integers and
$vv'\le4V_1^2$. Hence $|Y'|/Q\ge2$ and
$\vartheta\in[0,1)$, exactly the shifted-Suzuki
range. If $Y'<0$, complex conjugation reduces to $|Y'|$
without changing the norm.

The source's cotangent amplitudes are nonnegative,
nondecreasing on the positive clipped progression and at most
one. Their product has the same properties. Its finite
variation telescopes to \eqref{eq:suzuki-amplitude-variation}.
Apply Abel summation to the very same $A_t$, not to a constant
replacement; bound each partial interval by Suzuki and
multiply by $\mathcal V(A)\le2$. This proves
\eqref{eq:suzuki-clipped-conclusion}. In the formal source these
moves are \texttt{QuillonRLemma.suzuki\_AP},
\texttt{QuillonRMode.class\_amp\_bound},
\texttt{QuillonCot.ampC\_abel\_factor} and
\texttt{QuillonRPiece.R\_offdiag\_class}, which treats
$M\le N$ separately as the empty sum.
\end{proof}

\begin{remark}[Four channels and the two physical thresholds]
\label{rem:suzuki-four-channels}
The \emph{outer Fourier coefficients} are not the varying
$A_t$ in \cref{lem:clipped-suzuki}. The actual main and class
rows, for each $k\ge1$, are
\begin{center}
\small
\begin{tabular}{@{}lll@{}}
\toprule
Channel & Physical phase parameter $Y_k$ & Outer coefficient\\
\midrule
Main $+$ & $kH/r$ & $-2\mu\,\beta_k$\\
Main $-$ & $-kH/r$ & $+2\mu\,\beta_k$\\
Class $+$ & $kH/(6qr)$ & $D_q\beta_k$\\
Class $-$ & $-kH/(6qr)$ & $-D_q\beta_k$\\
\bottomrule
\end{tabular}
\end{center}
Here $D_q=2(6q)(48q)+1$, $|\mu|\le4$ and
$|\beta_k|\le1/k$, so the source proves the common bound
\begin{equation}\label{eq:suzuki-mode-coeff}
 |\mathrm{coef}_k|\le\frac{600q^2+8}{k}
\end{equation}
for every one of the four calls. The residue modulus is
$Q=6qr$ \emph{for all four rows}; the main and class phases
still have their different displayed denominators.
For each physical channel the source checks
$|Y_k|\ge8QV_1^2$, the actual clipped geometry, and
summability of \eqref{eq:suzuki-mode-coeff}. These
additional proofs occur in \texttt{R\_final} and
\texttt{R\_final\_L}, not in Abel's identity alone.
CORE $R$ uses $\log r\ge5\Lam^{\sigma_0}$;
LONG-right $RL$ proves its own $\log r\ge4\Lam^{\sigma_0}$.
The resulting distinct $125c/18$ and $64c/18$ constants
appear in \eqref{eq:vk-lower}. Neither the boxed clipped
lemma nor the exponent algebra by itself proves the later
fixed-$B$ bound \eqref{eq:r-fixed-B}; the full source also
pays packet extraction, Cauchy, Fourier coefficients and
exceptional/diagonal terms.
\end{remark}

\subsection{Exactly where CORE's 5 and LONG's 4 diverge}

The source's eventual parameter proofs reduce the oscillatory
part to $\exp\{-cL^3/(2L_m^2)\}$ with
$0<L_m\le3\Lam$. CORE R has the admissible lower bound
$L\ge5\Lam^{\sigma_0-1/200}$; LONG-right RL has only
$L\ge4\Lam^{\sigma_0-1/200}$. In either case let
$k\in\{4,5\}$. Then
\begin{equation}\label{eq:vk-lower}
 \frac{cL^3}{2L_m^2}
 \ge\frac{ck^3}{18}\,
 \Lam^{3\sigma_0-3/200-2}
 =\frac{ck^3}{18}\Lam^{3/200},
\end{equation}
using exactly
\begin{equation}\label{eq:positive-margin}
 3\left(\frac23+\frac1{100}\right)-\frac3{200}-2
 =\frac3{200}>0.
\end{equation}
The lower constants are therefore $125c/18$ for CORE
and $64c/18$ for LONG. For every $A>0$ and $d>0$,
$\exp(-d\Lam^{3/200})\le\Lam^{-A}$ eventually,
since $\Lam^{3/200}/\log\Lam\to\infty$.
This proves the \emph{parameter implication} that both
source settings can eventually pay $\Lam^{-200}$;
it does not by itself re-prove the full source's $R$-cell
estimate.

The separately supplied fixed-$B$ source theorems
\texttt{R\_final} and \texttt{R\_final\_L}, after the
packet extraction, four Fourier-mode estimates, coefficient
variation and source convexity, give
\begin{equation}\label{eq:r-fixed-B}
 \left|\sum_{C\text{ in physical }R\text{ or }RL\text{ mask}}
     \omegaR(C)\F_{q,a,D}(H,r,s)\right|
 \le4\Lam^{-90}.
\end{equation}
The same outer mass bound \eqref{eq:outer-mass} for the
fixed-$B$ sum yields
$4e^{10}5^{10}\Lam^{-80}$, absorbed into
$\Lam^{-79}\le\Lam^{-75}$ at sufficiently large $\Lam$.
The two physical results are consequently
$|R_d|,|RL_d|\le\Lam^{-75}$.
Crucially, \emph{$R\_final\_L$ is an independent LONG-specific
quantitative port}; it is not an application of CORE's
$R\_final$ after dropping its $5$-threshold hypothesis.

\section{The B-boundary centering: a finite numerator, an infinite subtraction}
\label{sec:centering}

This mechanism is an identity, not a logarithmic estimate.
Fix a first anchor representative $r\pmod q$, let
$I_q(r;h)=\{0<t<h:(r+t,q)=1\}$ and
$N_q(r;h)=|I_q(r;h)|$. Define the exact local exceptional-prime
weight
\begin{equation}\label{eq:low-weight}
 w_{q,r}(t)=\one_{(r+t,q)=1}
 \prod_{\substack{p\in\{2,3\}\\p\nmid q}}
 \frac{1-\nu_p(\{0,t\})/p}{(1-1/p)^2}.
\end{equation}
For high primes put
\begin{equation}\label{eq:eta-normalization}
 \eta_q(e)=\mu^2(e)\one_{(e,6q)=1}
    \prod_{p\mid e}\frac1{p-2},\quad
 C_{f,q}=\prod_{\substack{p\ge5\\p\nmid q}}
             \frac{p(p-2)}{(p-1)^2},\quad
 C_{f,q}\sum_{e\ge1}\frac{\eta_q(e)}e=1.
\end{equation}
The normalization is an equality of a convergent product
and an absolutely convergent reciprocal series. Define
\[
 R_e^w(q,r;h)=\sum_{\substack{0<t<h\\e\mid t}}w_{q,r}(t),
 \qquad S_1(q,r;h)=\sum_{t\in I_q(r;h)}
           \bigl(\SSer_q(\{0,t\})-1\bigr).
\]
The sealed Euler-bridge lemmas supply the exact expansion of
its positive part and the nonvanishing/convergence facts
needed to combine it with \eqref{eq:eta-normalization}.

\begin{proposition}[Why centering cannot be truncated]
\label{prop:centering-tail}
Under these exact normalization and summability identities,
\begin{equation}\label{eq:centering-identity}
 \boxed{S_1(q,r;h)=C_{f,q}\sum_{e\ge1}\eta_q(e)
     \left(R_e^w(q,r;h)-\frac{N_q(r;h)}e\right).}
\end{equation}
The positive term $R_e^w(q,r;h)$ is zero for $e\ge h$,
whereas the negative centering term generally is not.
If $N_q(r;h)>0$, truncating both pieces at $e<h$ omits
exactly the negative correction
\begin{equation}\label{eq:missing-negative-tail}
 -C_{f,q}N_q(r;h)\sum_{e\ge h}\frac{\eta_q(e)}e.
\end{equation}
\end{proposition}
\begin{proof}
The primewise free-pair selector identity and finite
positive divisor swap yield
$S_1=C_{f,q}\sum_e\eta_q(e)R_e^w-N_q$.
The full normalization in \eqref{eq:eta-normalization}
allows the \emph{exact} substitution
$N_q=C_{f,q}N_q\sum_{e\ge1}\eta_q(e)/e$.
Absolute summability permits termwise subtraction, proving
\eqref{eq:centering-identity}. For $e\ge h$, no integer
$0<t<h$ is divisible by $e$, so the positive term is zero.
The negative tail remains as \eqref{eq:missing-negative-tail},
and does not become zero through that support argument.
\end{proof}

At $q=1,h=2$ there is just one possible interior $t=1$.
Its active $p=2$ factor vanishes in \eqref{eq:low-weight},
so $R_e^w=0$ for every $e$, while $N_1(0;2)=1$.
By \eqref{eq:eta-normalization}, \eqref{eq:centering-identity}
gives $S_1(1,0;2)=-1$. Omitting the exceptional-prime
weight would produce the wrong positive progression mass.

For the second endpoint one \emph{must} reflect $t\mapsto h-t$
and use $r_b\equiv-b\pmod q$. The congruence
$a+h\equiv b\pmod q$ converts the second admissibility
condition into $(r_b+t,q)=1$. Thus the full printed B
quantity is $S_1(q,a;h)+S_1(q,r_b;h)$, with separate
low-prime weights and ordinary counts in both centered
series. This is the precise content of the formal reflected
Euler bridge, not a symmetry assumption imposed after
centering.

\section{The remaining physical cells: precise method contracts}
\label{sec:contracts}

The four demonstrations above treat five of the twelve physical
cells in detail: Band, EdgeL, WRL, R and RL. Seven additional
cells must be accounted for in the exact split; there is no
need to force each into a chapter. Write $\mathsf M_0$
for the inherited CORE predicate, $\mathsf M_L$ for
LONG, and $\mathsf W$ for the original CORE B-prime
packet. The tests
\[
 \mathsf G_C:\ s^{1+\epsp}\le r,
 \qquad \mathsf G_R:\ r^{1+3\epsp}\le s
\]
have the displayed equality conventions. Every item below
inherits the original decorated source, $B,C\ne\varnothing$,
$B\cap C=\varnothing$, cardinality window, signed row,
source prime exclusions and Central gate. The indicated bound
holds for the \emph{complete selected sum at one decoration},
eventually and uniformly in the remaining legal variables.

\subsection{CORE: three oscillatory regimes and a rough residual}

\paragraph{$W$ (first B packet; $\Lam^{-80}$).}
The additional mask is $\mathsf M_0\wedge\mathsf W$.
An eligible B-prime lies in the source packet
$[X_p(h),Y_p(\delta,s)]$ with the true active support and
coprimality. Ramar\'e extraction makes that prime the long
variable; the source then applies a fourth-moment/CRT
complete-sum estimate of square-root-of-fourth-moment
strength, with the required four physical Fourier rows.
This CORE input uses the source's $s^{3/4}$ modulus
power; it is \emph{not} the general-modulus Weil input
subsequently required by $WRL$.
The fixed-factor saving and outer cost are $-95+10=-85$,
then a constant absorbs five powers to reach $-80$.
The leaf is \path{GoldC/QuillonWPhysical.lean}.

\paragraph{$C$ (no B packet; $\Lam^{-15}$).}
Its mask is $\mathsf M_0\wedge\neg\mathsf W
\wedge\mathsf G_C$. The long first factor supplies a
range in which Parseval and Vaughan's identity convert
the relevant source Fourier sum to a suitable prime sum;
the exceptional portion is paid by Rankin's inequality.
The verified fixed-factor power $-30$ incurs the
specified outer $+10$ cost and absorbs a fixed coefficient
before the physical target $-15$.
See \path{GoldC/QuillonCPhysical.lean}.

\paragraph{$\mathrm{ResA}$ (no usable packet; $\Lam^{1/10}$).}
The mask is $\mathsf M_0\wedge\neg\mathsf W
\wedge\neg\mathsf G_C\wedge\mathsf G_R$
with $C\cap\JR(h,T,B)=\varnothing$. The failure of both
packet tests forces the two factors into distinct roughness
classes. The source bounds their \emph{weighted masses},
not the oscillation of a missing long variable. Its raw
power is
\[
 \frac1{4000}+\frac1{4000}+\frac1{40}
   =\frac{51}{2000}<\frac1{10},
\]
leaving $149/2000$ for constant absorption. There is no
additional physical outer $+10$ charge here because the
mass bound already covers the masked double sum.
See \path{GoldC/QuillonResPhysical.lean}.

The other two CORE cells, $R$ and Band, are treated in
\cref{sec:reciprocal,sec:band}. Their complements make
the five-way identity exhaustive; none of the above
methods is asserted outside its assigned mask.

\subsection{LONG-left: packet or Parseval}

\paragraph{$WL$ (left B packet; $\Lam^{-80}$).}
The mask retains LONG-left $s^{1+\delta}<r$ and the
eligible source W-packet in B. The extraction and
four-mode Fourier control are analogous to CORE $W$,
\emph{but the LONG-specific intervals and lower bounds
are reverified}; it is not a use of $W$ after discarding
its CORE hypothesis. The fixed-factor $-95$, outer $+10$
and fixed-constant absorption give $-80$.
See \path{GoldC/QuillonLW.lean}.

\paragraph{$CL$ (left with no packet; $\Lam^{-15}$).}
Its mask is LONG-left with the complementary
$\neg\mathsf W_L$. A dedicated long-range
Parseval/Vaughan decomposition and Rankin remainder
prove its version of $-30+10$, then absorb constants to
$-15$. The source's ranges are different from CORE's.
See \path{GoldC/QuillonLC.lean}.

\subsection{LONG-right: factors too numerous or packets missing}

\paragraph{$\mathrm{ManyL}$ (many C factors; $\Lam^{-50}$).}
After excluding the reciprocal packet
$C\cap\JR=\varnothing$, the criterion
$|C|\ge K_0(h)=\lceil100\log\Lam\rceil$
selects a Rankin factor-count estimate. Its factor
suppression pays $-100$; the cubic moment costs $+30$,
and the physical outer sum costs $+10$, leaving $-60$
before constant absorption to $-50$. This is not a
Kloosterman bound and does not require a W-prime packet.
See \path{GoldC/QuillonLRP.lean}.

\paragraph{$\mathrm{ResL}$ (off edge without W packet;
$\Lam^{7/10}$).}
The inherited tests are LONG-right, $C\cap\JR=\varnothing$,
$|C|<K_0(h)$ and $r\ge s^{1/2+3\eta}$; the decisive
additional test is $\mathcal W_L(B,C)=\varnothing$.
The two factors are rough enough to count directly.
The fixed-decoration raw power is
\[
 \frac1{2000}+\frac12+\frac1{40}
  =\frac{1051}{2000}<\frac7{10},
\]
leaving $349/2000$ of absorption slack. The half-power
is a genuine $(\eta)^{-5}=\Lam^{1/2}$ roughness
expense, not the edge's geometric width. This
physical mass already includes its outer cost.
See \path{GoldC/QuillonLRP.lean}.

The other LONG-right cases are $RL$ (the first reciprocal
packet), EdgeL (the thin square-root strip) and $WRL$
(the eligible off-edge B packet). They were derived in
\cref{sec:reciprocal,sec:edge}. Every right-oriented
parent is routed into exactly one of these five regions.

\section{How the twelve cells pay for the theorem}\label{sec:accounting}

For the leaf estimates supplied in the sealed source, the
quantitative book has two levels. \emph{Fixed-factor inputs}
may require outer source summation; the positive mass bounds
already include their corresponding physical sums. This is
why adding $+10$ to every row of the exponent table would
be incorrect.

\begin{center}\small
\begin{tabular}{@{}lrrr@{}}
\toprule
Cell & Before constants & Physical target & Unused exponent\\\midrule
$W,WL,WRL$ & $-85$ & $-80$ & $5$\\
$C,CL$ & $-20$ & $-15$ & $5$\\
$R,RL$ & $-80$ & $-75$ & $5$\\
$\mathrm{ManyL}$ & $-60$ & $-50$ & $10$\\
$\mathrm{ResA}$ & $51/2000$ & $1/10$ & $149/2000$\\
$\mathrm{Band}$ & $1991/2000$ & $499/500$ & $1/400$\\
$\mathrm{EdgeL}$ & $3741/4000$ & $19/20$ & $59/4000$\\
$\mathrm{ResL}$ & $1051/2000$ & $7/10$ & $349/2000$\\
\bottomrule
\end{tabular}
\end{center}

For every fixed decoration, exact disjointness followed by
the triangle inequality gives
\begin{equation}\label{eq:per-decoration-conclusion}
 |\mathcal C_{\csmall,d}|\le5\Lam^{499/500},
 \qquad |\mathcal C_{\clong,d}|\le7\Lam^{19/20}.
\end{equation}
A difference of signs is preserved up to this step; one
never multiplies the five or seven exponent losses.

The decoration masses are nonnegative. Their separate
supplier bound has the form
\begin{equation}\label{eq:decor-bound}
 \sum_dV_h(d)\le e^{10}\left(1+\frac{\log\Lam}{\log2}\right)^8
 \le256e^{10}\Lam^{1/1000}
\end{equation}
eventually. Multiplying \emph{once} by this aggregate bound,
\begin{equation}\label{eq:final-cost}
 \boxed{
 |\mathcal C_{\csmall}|\le1280e^{10}\Lam^{999/1000},\quad
 |\mathcal C_{\clong}|\le1792e^{10}\Lam^{951/1000}.}
\end{equation}
Both exponents are strictly smaller than one. The source's
positive fixed-$q$ HIGH contribution is $O_q(1)$ and the
zero-cardinality LOW sector with $J_{\rm low}\equiv0$
is \emph{exactly} zero. Every original parent not in HIGH
is then in MIDDLE and is covered by CORE or LONG.

The folded-to-active Central difference retains a separate
inverse-row and density-row correction, bounded by
\[
 (6q+2)\,4(1+\Lam)^3Z(h)^{-1/2}e^{K_0(h)},
\]
which tends to zero in the proved eventual range. The
independently paid positive Vaaler and shifted Suzuki
premises then control the \emph{other} regions in the
four-way quarter-middle decomposition, including Upper;
the four regions are requested at tolerance
$\varepsilon/4$. The source's far-quarter transport
and source-to-kernel normalization use their own
$\varepsilon/2+\varepsilon/2$ allocation. Those splits
change the chosen tolerances, \emph{not} any physical
logarithmic exponent. Finally the exact printed Euler
identification transports the native coefficient difference
to the literal statement
\[
 A_{1,q}(a,b;h)-B_{1,q}(a,b;h)=o_q(\log h).
\]
The height threshold depends on the fixed modulus and the
requested error but is selected \emph{before} the endpoint
classes. The method is asymptotic: the constant-absorption
thresholds are not claimed to be numerically practical.

\section{Source waypoints and independent obligations}
\label{sec:sources}

The sealed October 2026 R0003 development is the source for the
imported analytic lemmas. The separate \emph{source-obligation ledger}
provides exact archive and file hashes, declaration names, line ranges,
assumptions and falsification questions. It distinguishes what has been
derived in this supplement from what is supplied by the formal source;
it is not presented as a compiler-extracted proof-term graph.

\paragraph{Exact Euler identities.}
The F-01 printed-literal bridge supplies the centered $S_1$ identity,
the convergence of its series and the reflected-anchor normalization.
Those are equalities, not estimates.

\paragraph{Geometric mass and analytic cancellation.}
For Band, the named inputs are \texttt{interval\_mass},
\texttt{rough\_ev} and \texttt{Band\_physical}.
For the edge, they are \texttt{rough\_ev4} and
\texttt{EdgeL\_physical}; off the edge, the paid general
Kloosterman estimate, \texttt{WeilHyp}, the coupled gcd moment and
\texttt{W\_final\_LR} are required together. The shifted Suzuki
input enters the two distinct physical estimates
\texttt{R\_final} and \texttt{R\_final\_L}; neither supplies the
other's source hypotheses for free.

\paragraph{Aggregation and remaining review.}
The exact five- and seven-cell identities and their decorated
sublogarithmic estimates precede the zero-LOW specialization and
literal printed terminal. A cold analyst should independently attack
the roughness and harmonic-interval suppliers, the Weil cofactor/gcd
contract, and the clipped-coefficient Abel/Suzuki interface. Those
are real mathematical review tasks. The deep inputs are established
in the sealed formal development, but they have \emph{not} all been
independently re-proved by handwritten arguments here. The external
Lean replay succeeded; handwritten mathematical audit remains open.


\begingroup
\small
\begin{thebibliography}{12}
\setlength{\itemsep}{1pt}

\bibitem{MontgomerySoundararajan2004}
H.~L.~Montgomery and K.~Soundararajan,
\emph{Primes in short intervals},
Commun. Math. Phys. \textbf{252} (2004), 589--617.
DOI: 10.1007/s00220-004-1222-4.

\bibitem{Vaaler1985}
J.~D.~Vaaler,
\emph{Some extremal functions in Fourier analysis},
Bull. Amer. Math. Soc. (N.S.) \textbf{12} (1985), no.~2, 183--216.

\bibitem{Weil1948}
A.~Weil,
\emph{On some exponential sums},
Proc. Natl. Acad. Sci. USA \textbf{34} (1948), 204--207.

\bibitem{Estermann1961}
T.~Estermann,
\emph{On Kloosterman's sum},
Mathematika \textbf{8} (1961), no.~1, 83--86.

\bibitem{IwaniecKowalski2004}
H.~Iwaniec and E.~Kowalski,
\emph{Analytic Number Theory},
Amer. Math. Soc. Colloquium Publications, vol.~53,
American Mathematical Society, Providence, 2004.

\bibitem{Karatsuba1971}
A.~A.~Karatsuba,
\emph{Estimates for trigonometric sums by Vinogradov's method,
and some applications},
Proc. Steklov Inst. Math. \textbf{112} (1971), 251--265
(English translation, 1973).

\bibitem{Ramare2010}
O.~Ramar\'e,
\emph{On Bombieri's asymptotic sieve},
J. Number Theory \textbf{130} (2010), no.~5, 1155--1189.

\bibitem{Vaughan1980}
R.~C.~Vaughan,
\emph{An elementary method in prime number theory},
Acta Arith. \textbf{37} (1980), 111--115.

\bibitem{GoldstonLedoan2011}
D.~A.~Goldston and A.~H.~Ledoan,
\emph{On the differences between consecutive prime numbers, I},
arXiv:1111.3380 (2011).

\bibitem{LemkeOliverSound2016}
R.~J.~Lemke Oliver and K.~Soundararajan,
\emph{Unexpected biases in the distribution of consecutive primes},
Proc. Natl. Acad. Sci. USA \textbf{113} (2016),
no.~31, E4446--E4454.

\bibitem{Kuperberg2025}
V.~Kuperberg,
\emph{Sums of singular series along arithmetic progressions
and with smooth weights},
Int. J. Number Theory \textbf{21} (2025), no.~1, 53--74.

\end{thebibliography}
\endgroup

\end{document}
